\documentclass{book}
\usepackage[utf8]{inputenc}
\usepackage[T1]{fontenc}
\usepackage{lmodern}
\usepackage[spanish]{babel}
\addto\captionsspanish{\renewcommand{\contentsname}{Índice}}
\AtBeginDocument{\spanishdeactivate{"~<>}}
\addto\extrasspanish{\spanishdeactivate{"~<>}}
\usepackage[a4paper,margin=2.5cm]{geometry}
\usepackage{titlesec}
\titleformat{\chapter}[display]
  {\normalfont\filright}{\large\scshape\chaptertitlename\ \thechapter}{1ex}{\huge\bfseries}
\titlespacing*{\chapter}{0pt}{2.5ex plus 1ex minus .2ex}{2.3ex}
\titleformat{name=\chapter,numberless}[display]
  {\normalfont\filright}{}{0pt}{\huge\bfseries}
\titleformat{\section}{\Large\bfseries}{{\thesection}}{0.8em}{}
\titleformat{\subsection}{\large\bfseries}{{\thesubsection}}{0.7em}{}
\setcounter{secnumdepth}{3}
\setcounter{tocdepth}{2}
\usepackage{graphicx}
\usepackage{tikz}
\usetikzlibrary{arrows.meta,positioning,calc,shapes.geometric}
\usepackage{xcolor}
\usepackage{amsmath}
\usepackage{amssymb}
\usepackage{array}
\usepackage{longtable}
\newcounter{studiumfigure}
\definecolor{studiumInk}{RGB}{28,40,58}
\definecolor{studiumRule}{RGB}{28,40,58}
\definecolor{studiumTipBack}{RGB}{232,244,236}
\definecolor{studiumTipFrame}{RGB}{27,107,58}
\definecolor{studiumDefBack}{RGB}{232,240,250}
\definecolor{studiumDefFrame}{RGB}{24,74,130}
\definecolor{studiumWorkedBack}{RGB}{255,244,230}
\definecolor{studiumWorkedFrame}{RGB}{138,78,12}
\definecolor{studiumCheckBack}{RGB}{243,236,248}
\definecolor{studiumCheckFrame}{RGB}{96,48,130}
\definecolor{studiumCodeBack}{RGB}{245,245,242}
\definecolor{studiumCodeRule}{RGB}{48,48,52}
\usepackage{listings}
\lstdefinelanguage{Rust}{
  morekeywords={as,async,await,break,const,continue,crate,dyn,else,enum,extern,false,fn,for,if,impl,in,let,loop,match,mod,move,mut,pub,ref,return,self,Self,static,struct,super,trait,true,type,unsafe,use,where,while},
  sensitive=true,
  morecomment=[l]{//},
  morecomment=[s]{/*}{*/},
  morestring=[b]{"}
}
\lstdefinelanguage{Go}{
  morekeywords={break,case,chan,const,continue,default,defer,else,fallthrough,for,func,go,goto,if,import,interface,map,package,range,return,select,struct,switch,type,var},
  sensitive=true,
  morecomment=[l]{//},
  morecomment=[s]{/*}{*/},
  morestring=[b]{"}
}
\lstdefinelanguage{JavaScript}{
  morekeywords={break,case,catch,class,const,continue,debugger,default,delete,do,else,export,extends,finally,for,function,if,import,in,instanceof,let,new,return,super,switch,this,throw,try,typeof,var,void,while,with,yield},
  sensitive=true,
  morecomment=[l]{//},
  morecomment=[s]{/*}{*/},
  morestring=[b]{"},
  morestring=[b]{'},
}
\lstdefinestyle{studium}{
  basicstyle=\ttfamily\small,
  columns=fullflexible,
  keepspaces=true,
  showstringspaces=false,
  breaklines=true,
  breakatwhitespace=false,
  frame=leftline,
  framerule=1.15pt,
  rulecolor=\color{studiumCodeRule},
  backgroundcolor=\color{studiumCodeBack},
  xleftmargin=1.15em,
  framexleftmargin=0.75em,
  aboveskip=0.9em,
  belowskip=0.7em,
  tabsize=4,
  keywordstyle=\ttfamily\bfseries,
  commentstyle=\ttfamily\itshape,
  stringstyle=\ttfamily
}
\lstset{style=studium}
\usepackage[breakable,skins]{tcolorbox}
\tcbset{
  studiumtip/.style={colback=studiumTipBack,colframe=studiumTipFrame,colbacktitle=studiumTipFrame,coltitle=white,coltext=black,fonttitle=\bfseries,boxrule=0.65pt,arc=0.7pt,left=1.7mm,right=1.7mm,top=1.3mm,bottom=1.3mm,before skip=10pt,after skip=10pt},
  studiumdef/.style={colback=studiumDefBack,colframe=studiumDefFrame,colbacktitle=studiumDefFrame,coltitle=white,coltext=black,fonttitle=\bfseries,boxrule=0.65pt,arc=0.7pt,left=1.7mm,right=1.7mm,top=1.3mm,bottom=1.3mm,before skip=10pt,after skip=10pt},
  studiumworked/.style={colback=studiumWorkedBack,colframe=studiumWorkedFrame,colbacktitle=studiumWorkedFrame,coltitle=white,coltext=black,fonttitle=\bfseries,boxrule=0.65pt,arc=0.7pt,left=1.7mm,right=1.7mm,top=1.3mm,bottom=1.3mm,before skip=10pt,after skip=10pt},
  studiumcheck/.style={colback=studiumCheckBack,colframe=studiumCheckFrame,colbacktitle=studiumCheckFrame,coltitle=white,coltext=black,fonttitle=\bfseries,boxrule=0.65pt,arc=0.7pt,left=1.7mm,right=1.7mm,top=1.3mm,bottom=1.3mm,before skip=10pt,after skip=10pt}
}
\titleformat{\section}{\Large\bfseries}{\thesection}{0.8em}{}
\titleformat{\subsection}{\large\bfseries}{\thesubsection}{0.7em}{}
\setcounter{secnumdepth}{3}
\setcounter{tocdepth}{2}
\usepackage{fancyhdr}
\pagestyle{fancy}
\fancyhf{}
\fancyhead[LE]{\small\thepage\hspace{1.5em}\scshape\nouppercase{\leftmark}}
\fancyhead[RO]{\small\scshape\nouppercase{\rightmark}\hspace{1.5em}\thepage}
\fancyfoot[C]{\small Borrador}
\renewcommand{\headrulewidth}{0.4pt}
\renewcommand{\footrulewidth}{0pt}
\fancypagestyle{plain}{%
  \fancyhf{}
  \fancyfoot[C]{\small Borrador}
  \renewcommand{\headrulewidth}{0pt}
  \renewcommand{\footrulewidth}{0pt}
}
\usepackage[hyphens]{url}
\usepackage[hidelinks]{hyperref}
\setlength{\emergencystretch}{1.5em}
\makeatletter
\@openrightfalse
\let\cleardoublepage\clearpage
\makeatother
\begin{document}
\frontmatter
\thispagestyle{empty}
\setlength{\parindent}{0pt}
\vspace*{0.08\textheight}
{\normalsize\scshape\color{studiumInk} Studium\par}
{\small Enciclopedia de Programación e Ingeniería del Software\par}
\vspace{0.55em}
{\color{studiumRule}\rule{\linewidth}{0.9pt}\par}
\vspace{0.16\textheight}
{\raggedright\Huge\bfseries Volumen IV: Monografías de Lenguajes de Sistemas \(-\) C, C++, Zig, Ensamblador, Ada y Fortran\par}
\vspace{0.9em}
{\raggedright\Large\color{studiumInk} Parte V (Bloques A\(-\)H): Memoria manual, RAII, metaprogramación con Concepts, ABI de bajo nivel y cómputo científico\par}
\vfill
{\color{studiumRule}\rule{0.36\linewidth}{0.45pt}\par}
\vspace{1.05em}
{\large\scshape Borrador\par}
\vspace{0.45em}
{\normalsize Edición 0.1.0 --- No publicado\par}
\vspace{0.4em}
{\normalsize\today\par}
\vspace*{0.08\textheight}
\clearpage
\chapter{Prefacio}
\noindent Este libro presenta el tema con el lenguaje del curso. Cada capítulo abre con una entrada breve, sigue con la explicación y cierra con un problema resuelto y una autoficha.
\chapter{Cómo usar este libro}
\noindent Lee el capítulo seguido. El consejo, la definición, el problema resuelto y la autoficha van en recuadros. La explicación es el cuerpo del texto.

\noindent Un capítulo lleva una entrada en cursiva, la explicación en el cuerpo, como mucho un consejo, definiciones solo al introducir un término, un problema resuelto y una autoficha.
\tableofcontents
\mainmatter

\part{Monografía Autónoma de C (ISO/IEC 9899:2018 / C23)}

\chapter{Monografía de C (I \(-\) Bloques A, B y C): Historia, modelo de traducción ISO C17/C23, toolchain y nivel principiante}
\label{chap:cap01-c-fundamentos-modelo-compilacion-tipos-principiante}
% studium:record eyJpZCI6ICJQQVItMDAwMSIsICJzZWN0aW9uIjogImNhcDAxLWMtZnVuZGFtZW50b3MtbW9kZWxvLWNvbXBpbGFjaW9uLXRpcG9zLXByaW5jaXBpYW50ZSJ9

\noindent\textit{Este capítulo inicia la monografía autónoma del lenguaje C (Bloques A, B y C). C es la lingua franca de los sistemas operativos (Linux, Windows NT, XNU/macOS), los motores de bases de datos (PostgreSQL, SQLite, Redis), los intérpretes de lenguajes (CPython, Lua) y los sistemas empotrados. El lector aprenderá a programar en C moderno (C17/C23) desde cero con disciplina de tipos y herramientas de verificación rigurosas.}
% studium:end-record eyJpZCI6ICJQQVItMDAwMSIsICJzZWN0aW9uIjogImNhcDAxLWMtZnVuZGFtZW50b3MtbW9kZWxvLWNvbXBpbGFjaW9uLXRpcG9zLXByaW5jaXBpYW50ZSJ9
% studium:record eyJpZCI6ICJGSUctMDAwMSIsICJzZWN0aW9uIjogImNhcDAxLWMtZnVuZGFtZW50b3MtbW9kZWxvLWNvbXBpbGFjaW9uLXRpcG9zLXByaW5jaXBpYW50ZSJ9

\par\vspace{\baselineskip}
\noindent\begin{minipage}{\linewidth}
\centering
\begingroup
\setbox0=\hbox{\begin{tikzpicture}
\tikzset{node distance=1.6cm, >=stealth, thick, font=\small}
\node[draw=blue!70!black, fill=blue!8, rounded corners, minimum width=2.2cm, minimum height=0.9cm, align=center] (src) {\texttt{arena.c}\\+\,\texttt{arena.h}};
\node[draw=blue!70!black, fill=blue!8, rounded corners, minimum width=2.3cm, minimum height=0.9cm, align=center, right=0.5cm of src] (cpp) {Preprocesador\\(\texttt{cc -E} $\to$ \texttt{.i})};
\node[draw=orange!75!black, fill=orange!12, rounded corners, minimum width=2.3cm, minimum height=0.9cm, align=center, right=0.5cm of cpp] (cc1) {Compilador\\(\texttt{cc -S} $\to$ \texttt{.s})};
\node[draw=teal!70!black, fill=teal!10, rounded corners, minimum width=2.3cm, minimum height=0.9cm, align=center, right=0.5cm of cc1] (as) {Ensamblador\\(\texttt{as} $\to$ \texttt{.o})};
\node[draw=green!60!black, fill=green!10, rounded corners, minimum width=2.2cm, minimum height=0.9cm, align=center, right=0.5cm of as] (ld) {Enlazador\\(\texttt{ld} $\to$ Binario)};
\draw[->] (src) -- (cpp);
\draw[->] (cpp) -- (cc1);
\draw[->] (cc1) -- (as);
\draw[->] (as) -- (ld);
\end{tikzpicture}}%
\count0=\ht0
\advance\count0 by \dp0
\count2=\wd0
\count4=\dimexpr\textheight/100*38\relax
\count6=\linewidth
\ifnum\count0<1
\count0=1
\fi
\ifnum\count2<1
\count2=1
\fi
\divide\count0 by 4096
\divide\count2 by 4096
\divide\count4 by 4096
\divide\count6 by 4096
\ifnum\count0<1
\count0=1
\fi
\ifnum\count2<1
\count2=1
\fi
\ifnum\count4<1
\count4=1
\fi
\ifnum\count6<1
\count6=1
\fi
\count8=\count2
\ifnum\count0>\count4
\count8=\numexpr\count2*\count4/\count0\relax
\fi
\ifnum\count8>\count6
\resizebox{\linewidth}{!}{\begin{tikzpicture}
\tikzset{node distance=1.6cm, >=stealth, thick, font=\small}
\node[draw=blue!70!black, fill=blue!8, rounded corners, minimum width=2.2cm, minimum height=0.9cm, align=center] (src) {\texttt{arena.c}\\+\,\texttt{arena.h}};
\node[draw=blue!70!black, fill=blue!8, rounded corners, minimum width=2.3cm, minimum height=0.9cm, align=center, right=0.5cm of src] (cpp) {Preprocesador\\(\texttt{cc -E} $\to$ \texttt{.i})};
\node[draw=orange!75!black, fill=orange!12, rounded corners, minimum width=2.3cm, minimum height=0.9cm, align=center, right=0.5cm of cpp] (cc1) {Compilador\\(\texttt{cc -S} $\to$ \texttt{.s})};
\node[draw=teal!70!black, fill=teal!10, rounded corners, minimum width=2.3cm, minimum height=0.9cm, align=center, right=0.5cm of cc1] (as) {Ensamblador\\(\texttt{as} $\to$ \texttt{.o})};
\node[draw=green!60!black, fill=green!10, rounded corners, minimum width=2.2cm, minimum height=0.9cm, align=center, right=0.5cm of as] (ld) {Enlazador\\(\texttt{ld} $\to$ Binario)};
\draw[->] (src) -- (cpp);
\draw[->] (cpp) -- (cc1);
\draw[->] (cc1) -- (as);
\draw[->] (as) -- (ld);
\end{tikzpicture}}%
\else
\ifnum\count0>\count4
\resizebox{!}{\dimexpr\textheight/100*38\relax}{\begin{tikzpicture}
\tikzset{node distance=1.6cm, >=stealth, thick, font=\small}
\node[draw=blue!70!black, fill=blue!8, rounded corners, minimum width=2.2cm, minimum height=0.9cm, align=center] (src) {\texttt{arena.c}\\+\,\texttt{arena.h}};
\node[draw=blue!70!black, fill=blue!8, rounded corners, minimum width=2.3cm, minimum height=0.9cm, align=center, right=0.5cm of src] (cpp) {Preprocesador\\(\texttt{cc -E} $\to$ \texttt{.i})};
\node[draw=orange!75!black, fill=orange!12, rounded corners, minimum width=2.3cm, minimum height=0.9cm, align=center, right=0.5cm of cpp] (cc1) {Compilador\\(\texttt{cc -S} $\to$ \texttt{.s})};
\node[draw=teal!70!black, fill=teal!10, rounded corners, minimum width=2.3cm, minimum height=0.9cm, align=center, right=0.5cm of cc1] (as) {Ensamblador\\(\texttt{as} $\to$ \texttt{.o})};
\node[draw=green!60!black, fill=green!10, rounded corners, minimum width=2.2cm, minimum height=0.9cm, align=center, right=0.5cm of as] (ld) {Enlazador\\(\texttt{ld} $\to$ Binario)};
\draw[->] (src) -- (cpp);
\draw[->] (cpp) -- (cc1);
\draw[->] (cc1) -- (as);
\draw[->] (as) -- (ld);
\end{tikzpicture}}%
\else
\begin{tikzpicture}
\tikzset{node distance=1.6cm, >=stealth, thick, font=\small}
\node[draw=blue!70!black, fill=blue!8, rounded corners, minimum width=2.2cm, minimum height=0.9cm, align=center] (src) {\texttt{arena.c}\\+\,\texttt{arena.h}};
\node[draw=blue!70!black, fill=blue!8, rounded corners, minimum width=2.3cm, minimum height=0.9cm, align=center, right=0.5cm of src] (cpp) {Preprocesador\\(\texttt{cc -E} $\to$ \texttt{.i})};
\node[draw=orange!75!black, fill=orange!12, rounded corners, minimum width=2.3cm, minimum height=0.9cm, align=center, right=0.5cm of cpp] (cc1) {Compilador\\(\texttt{cc -S} $\to$ \texttt{.s})};
\node[draw=teal!70!black, fill=teal!10, rounded corners, minimum width=2.3cm, minimum height=0.9cm, align=center, right=0.5cm of cc1] (as) {Ensamblador\\(\texttt{as} $\to$ \texttt{.o})};
\node[draw=green!60!black, fill=green!10, rounded corners, minimum width=2.2cm, minimum height=0.9cm, align=center, right=0.5cm of as] (ld) {Enlazador\\(\texttt{ld} $\to$ Binario)};
\draw[->] (src) -- (cpp);
\draw[->] (cpp) -- (cc1);
\draw[->] (cc1) -- (as);
\draw[->] (as) -- (ld);
\end{tikzpicture}%
\fi
\fi
\endgroup
\par\nopagebreak
{\small \refstepcounter{studiumfigure}\label{fig:fig-c-translation-pipeline}\textbf{Figura~\thestudiumfigure.} Tubería clásica de compilación separada en C: el preprocesador expande directivas sobre cada unidad de traducción (.i), el compilador emite código ensamblador (.s), el ensamblador genera objetos reubicables (.o) y el enlazador resuelve símbolos externos.\par}
\end{minipage}
\par\vspace{\baselineskip}
% studium:end-record eyJpZCI6ICJGSUctMDAwMSIsICJzZWN0aW9uIjogImNhcDAxLWMtZnVuZGFtZW50b3MtbW9kZWxvLWNvbXBpbGFjaW9uLXRpcG9zLXByaW5jaXBpYW50ZSJ9

\begin{tcolorbox}[title={Consejo}, studiumtip, breakable]
% studium:record eyJpZCI6ICJQQVItMDAwNSIsICJzZWN0aW9uIjogImNhcDAxLWMtZnVuZGFtZW50b3MtbW9kZWxvLWNvbXBpbGFjaW9uLXRpcG9zLXByaW5jaXBpYW50ZSJ9

Incluye siempre guardas de inclusión (\texttt{\#ifndef MODULO\_H} / \texttt{\#define MODULO\_H} / \texttt{\#endif}) en todas tus cabeceras \texttt{.h} y nunca declares variables globales con definición en una cabecera sin \texttt{extern}.
% studium:end-record eyJpZCI6ICJQQVItMDAwNSIsICJzZWN0aW9uIjogImNhcDAxLWMtZnVuZGFtZW50b3MtbW9kZWxvLWNvbXBpbGFjaW9uLXRpcG9zLXByaW5jaXBpYW50ZSJ9
\end{tcolorbox}

\begin{tcolorbox}[title={Definición}, studiumdef, breakable]
% studium:record eyJpZCI6ICJQQVItMDAwNCIsICJzZWN0aW9uIjogImNhcDAxLWMtZnVuZGFtZW50b3MtbW9kZWxvLWNvbXBpbGFjaW9uLXRpcG9zLXByaW5jaXBpYW50ZSJ9

Comportamiento Indefinido (Undefined Behavior, ISO C17 §3.4.3): uso de una construcción no portable o de datos erróneos para los cuales el estándar internacional no impone requisito alguno, permitiendo que el compilador optimice bajo la premisa axiomática de que dicha situación es inalcanzable.
% studium:end-record eyJpZCI6ICJQQVItMDAwNCIsICJzZWN0aW9uIjogImNhcDAxLWMtZnVuZGFtZW50b3MtbW9kZWxvLWNvbXBpbGFjaW9uLXRpcG9zLXByaW5jaXBpYW50ZSJ9
\end{tcolorbox}

\begin{tcolorbox}[title={Autoficha}, studiumcheck, breakable]
% studium:record eyJpZCI6ICJQQVItMDAwNiIsICJzZWN0aW9uIjogImNhcDAxLWMtZnVuZGFtZW50b3MtbW9kZWxvLWNvbXBpbGFjaW9uLXRpcG9zLXByaW5jaXBpYW50ZSJ9

¿Por qué comparar un índice \texttt{int i = -1} con \texttt{size\_t n = 10} mediante \texttt{if (i < n)} evalúa a \texttt{false} en C debido a las conversiones aritméticas usuales?
% studium:end-record eyJpZCI6ICJQQVItMDAwNiIsICJzZWN0aW9uIjogImNhcDAxLWMtZnVuZGFtZW50b3MtbW9kZWxvLWNvbXBpbGFjaW9uLXRpcG9zLXByaW5jaXBpYW50ZSJ9
\end{tcolorbox}

\section{Bloques A y B: Historia de C, modelo de máquina abstracta, GCC/Clang, CMake y flags de diagnóstico}
% studium:record eyJpZCI6ICJQQVItMDAwMiIsICJzZWN0aW9uIjogImNhcDAxLWMtZnVuZGFtZW50b3MtbW9kZWxvLWNvbXBpbGFjaW9uLXRpcG9zLXByaW5jaXBpYW50ZSJ9

Creado por Dennis Ritchie entre 1969 y 1973 en los Laboratorios Bell para reescribir el núcleo de UNIX (originalmente en ensamblador PDP-11), C fue estandarizado por ANSI en 1989 (C89/C90) y por ISO en las revisiones C99 (que introdujo \texttt{<stdint.h>}, \texttt{<stdbool.h>}, inicializadores designados y el calificador \texttt{restrict}), C11 (modelo de memoria multihilo y \texttt{\_Atomic}), C17 y C23 (\texttt{nullptr}, \texttt{constexpr}, \texttt{typeof}, \texttt{<stdckdint.h>}). Como muestra la Figura \textbackslash{}ref\{fig:fig-c-translation-pipeline\}, cada fichero \texttt{.c} junto con las cabeceras \texttt{.h} incluidas por \texttt{\#include} constituye una *unidad de traducción* independiente. En ingeniería profesional, ningún programa en C debe compilarse sin activar las banderas estrictas de diagnóstico y los sanitizadores dinámicos: \texttt{clang -std=c17 -O2 -Wall -Wextra -Wpedantic -Werror -Wshadow -Wconversion -fsanitize=address,undefined}.
% studium:end-record eyJpZCI6ICJQQVItMDAwMiIsICJzZWN0aW9uIjogImNhcDAxLWMtZnVuZGFtZW50b3MtbW9kZWxvLWNvbXBpbGFjaW9uLXRpcG9zLXByaW5jaXBpYW50ZSJ9

\section{Bloque C: Tipos fundamentales, <stdint.h>, promociones aritméticas, operadores y control de flujo}
% studium:record eyJpZCI6ICJQQVItMDAwMyIsICJzZWN0aW9uIjogImNhcDAxLWMtZnVuZGFtZW50b3MtbW9kZWxvLWNvbXBpbGFjaW9uLXRpcG9zLXByaW5jaXBpYW50ZSJ9

Dado que el estándar C solo fija anchuras mínimas para \texttt{short}, \texttt{int}, \texttt{long} y \texttt{long long}, el código de sistemas portable utiliza siempre los tipos de anchura exacta de \texttt{<stdint.h>}: \texttt{uint8\_t}, \texttt{uint16\_t}, \texttt{uint32\_t}, \texttt{uint64\_t}, \texttt{int32\_t}, \texttt{int64\_t}, además de \texttt{size\_t} (para tamaños y conteos de objetos en memoria) y \texttt{ptrdiff\_t} (para diferencias entre punteros). Una regla crítica de la máquina abstracta de C es la asimetría entre enteros sin signo y con signo: mientras que los enteros sin signo (\texttt{uint32\_t}, \texttt{size\_t}) operan por definición en aritmética modular \(\bmod 2^w\), el desbordamiento de un entero **con signo** (\texttt{int32\_t}) constituye **comportamiento indefinido** (*Undefined Behavior*, UB), lo que autoriza al optimizador a asumir que jamás ocurre y a eliminar comprobaciones ingenuas a posteriori como \texttt{if (a + b < a)}. Por ello, el desbordamiento con signo debe prevenirse **antes** de ejecutar la suma:
\begin{lstlisting}[style=studium,language={C}]
#include <stdint.h>
#include <stdbool.h>

bool safe_add_i32(int32_t a, int32_t b, int32_t *out) {
    if ((b > 0 && a > INT32_MAX - b) || (b < 0 && a < INT32_MIN - b)) {
        return false;
    }
    *out = a + b;
    return true;
}
\end{lstlisting}
En C23, la cabecera estándar \texttt{<stdckdint.h>} estandariza esta operación mediante la primitiva intrínseca \texttt{ckd\_add(out, a, b)}.
% studium:end-record eyJpZCI6ICJQQVItMDAwMyIsICJzZWN0aW9uIjogImNhcDAxLWMtZnVuZGFtZW50b3MtbW9kZWxvLWNvbXBpbGFjaW9uLXRpcG9zLXByaW5jaXBpYW50ZSJ9

\begin{tcolorbox}[title={Problema resuelto}, studiumworked, breakable]
% studium:record eyJpZCI6ICJQUkItMDAwMSIsICJzZWN0aW9uIjogImNhcDAxLWMtZnVuZGFtZW50b3MtbW9kZWxvLWNvbXBpbGFjaW9uLXRpcG9zLXByaW5jaXBpYW50ZSJ9

\noindent\textbf{Enunciado}
En una arquitectura de 32 bits para enteros \texttt{uint32\_t} (\(w = 32\), módulo \(2^{32} = 4294967296\)), un entero con signo con valor \(-100\) se convierte implícitamente a \texttt{uint32\_t} durante una comparación. Calcula el valor numérico exacto que toma tras la conversión según la regla \(2^{32} - 100\).

\noindent\textbf{Datos}
\begin{itemize}
\item $w = 32\,\mathrm{bits}$ — Anchura en bits de uint\(32_t\)
\item $x = -100$ — Valor con signo original
\end{itemize}

\noindent\textbf{Incógnita}
Magnitud objetivo solicitada en el enunciado

\noindent\textbf{Modelo físico}

\noindent\textit{Conversión a entero sin signo módulo 2\textasciicircum{}32 (ISO C §6.3.1.3)}

\[
U = 2^{32} - 100
\]

\noindent\textbf{Hipótesis}
\begin{itemize}
\item Evaluación determinista bajo el modelo formal del capítulo
\end{itemize}

\noindent\textbf{Desarrollo}
\begin{enumerate}
\item
Sumamos \(2^32\) al valor negativo -100

\[
U = 2^{32} - 100
\]

\[
\left(\left(2\right)^{32}-100\right)
\]
\end{enumerate}

\noindent\textbf{Resultado}

\[
Resultado = 4294967196. El valor resultante en uint32_t es 4 294 967 196, razón por la cual `-Wsign-conversion` es indispensable.
\]

\noindent\textbf{Interpretación}
El valor resultante en uint\(32_t\) es 4 294 967 196, razón por la cual \texttt{-Wsign-conversion} es indispensable.

\noindent\textbf{Comprobación de los pasos}
\begin{itemize}
\item Paso 1: $4294967196$ — reproducido
\end{itemize}

\noindent\textbf{Comprobación de la respuesta} la aritmética del último paso reproduce la respuesta enunciada a la precisión con que está escrita. Es una reproducción aritmética; no valida el modelo físico ni sus hipótesis.
% studium:end-record eyJpZCI6ICJQUkItMDAwMSIsICJzZWN0aW9uIjogImNhcDAxLWMtZnVuZGFtZW50b3MtbW9kZWxvLWNvbXBpbGFjaW9uLXRpcG9zLXByaW5jaXBpYW50ZSJ9
\end{tcolorbox}

\begin{tcolorbox}[title={Ejercicios}, studiumworked, breakable]

\noindent\textbf{1 [PRB-0003] (Básico, Numérico)}
Calcula el valor en \texttt{uint16\_t} (\(w = 16\)) resultante de convertir \(-1\) a sin signo (\(2^{16} - 1\)).

\noindent\textbf{2 [PRB-0004] (Básico, Numérico)}
Un registro de control de 8 bits activa los bits 0, 3 y 7 (\$\(2^0 + 2^3 + 2^7\)\$). ¿Cuál es el valor decimal de la máscara \texttt{uint8\_t}?

\noindent\textbf{3 [PRB-0005] (Intermedio, Desarrollo por pasos)}
¿Cuántos bytes ocupa en total una estructura con 4 campos \texttt{uint32\_t} y 2 campos \texttt{uint64\_t} perfectamente alineados (\(4 \times 4 + 2 \times 8\))?

\noindent\textbf{4 [PRB-0006] (Intermedio, Diseño)}
Si compilamos un proyecto en C formado por 18 unidades de traducción \texttt{.c} y solo modificamos un fichero \texttt{.c}, ¿cuántas unidades de traducción evita recompilar \texttt{make}/\texttt{ninja}?
\end{tcolorbox}

\chapter{Monografía de C (II \(-\) Bloques D y E): Punteros, aritmética de direcciones, memoria dinámica, structs y strict aliasing}
\label{chap:cap02-c-punteros-memoria-dinamica-estructuras-aliasing}
% studium:record eyJpZCI6ICJQQVItMDAwNyIsICJzZWN0aW9uIjogImNhcDAyLWMtcHVudGVyb3MtbWVtb3JpYS1kaW5hbWljYS1lc3RydWN0dXJhcy1hbGlhc2luZyJ9

\noindent\textit{Este capítulo desarrolla los niveles intermedio y avanzado de C (Bloques D y E). Aquí desentrañamos la relación exacta entre arrays y punteros, la representación en memoria de estructuras y uniones, los riesgos de \texttt{malloc}/\texttt{free} individuales y por qué el patrón de ingeniería de sistemas por excelencia en C es el asignador por región (*Arena Allocator*).}
% studium:end-record eyJpZCI6ICJQQVItMDAwNyIsICJzZWN0aW9uIjogImNhcDAyLWMtcHVudGVyb3MtbWVtb3JpYS1kaW5hbWljYS1lc3RydWN0dXJhcy1hbGlhc2luZyJ9

\begin{tcolorbox}[title={Consejo}, studiumtip, breakable]
% studium:record eyJpZCI6ICJQQVItMDAxMSIsICJzZWN0aW9uIjogImNhcDAyLWMtcHVudGVyb3MtbWVtb3JpYS1kaW5hbWljYS1lc3RydWN0dXJhcy1hbGlhc2luZyJ9

Para reinterpretar los bits de un \texttt{float} como \texttt{uint32\_t} o leer enteros de un paquete de red en C, utiliza siempre \texttt{memcpy(\&destino, \&origen, sizeof(destino))}: es 100 \% seguro frente a strict aliasing y alineamiento, y el compilador lo optimiza a una sola instrucción de registro.
% studium:end-record eyJpZCI6ICJQQVItMDAxMSIsICJzZWN0aW9uIjogImNhcDAyLWMtcHVudGVyb3MtbWVtb3JpYS1kaW5hbWljYS1lc3RydWN0dXJhcy1hbGlhc2luZyJ9
\end{tcolorbox}

\begin{tcolorbox}[title={Definición}, studiumdef, breakable]
% studium:record eyJpZCI6ICJQQVItMDAxMCIsICJzZWN0aW9uIjogImNhcDAyLWMtcHVudGVyb3MtbWVtb3JpYS1kaW5hbWljYS1lc3RydWN0dXJhcy1hbGlhc2luZyJ9

Regla de Alias Estricto (Strict Aliasing, ISO C17 §6.5): restricción semántica que prohíbe acceder a un valor almacenado mediante un puntero de un tipo distinto no relacionado (salvo tipos de carácter \texttt{unsigned char}), permitiendo al compilador mantener valores en registros sin recargarlos desde memoria tras escrituras a otros tipos.
% studium:end-record eyJpZCI6ICJQQVItMDAxMCIsICJzZWN0aW9uIjogImNhcDAyLWMtcHVudGVyb3MtbWVtb3JpYS1kaW5hbWljYS1lc3RydWN0dXJhcy1hbGlhc2luZyJ9
\end{tcolorbox}

\begin{tcolorbox}[title={Autoficha}, studiumcheck, breakable]
% studium:record eyJpZCI6ICJQQVItMDAxMiIsICJzZWN0aW9uIjogImNhcDAyLWMtcHVudGVyb3MtbWVtb3JpYS1kaW5hbWljYS1lc3RydWN0dXJhcy1hbGlhc2luZyJ9

¿Por qué tras ejecutar \texttt{free(ptr)} en C es una buena práctica defensiva asignar \texttt{ptr = NULL} si el puntero permanece aún en ámbito?
% studium:end-record eyJpZCI6ICJQQVItMDAxMiIsICJzZWN0aW9uIjogImNhcDAyLWMtcHVudGVyb3MtbWVtb3JpYS1kaW5hbWljYS1lc3RydWN0dXJhcy1hbGlhc2luZyJ9
\end{tcolorbox}

\section{Bloque D: Punteros, arrays, cadenas terminadas en NUL, structs, unions y memoria dinámica}
% studium:record eyJpZCI6ICJQQVItMDAwOCIsICJzZWN0aW9uIjogImNhcDAyLWMtcHVudGVyb3MtbWVtb3JpYS1kaW5hbWljYS1lc3RydWN0dXJhcy1hbGlhc2luZyJ9

En C, un puntero \texttt{T *p} almacena la dirección virtual de un objeto de tipo \texttt{T}. La aritmética de punteros está escalada automáticamente por el tamaño del tipo apuntado: si \texttt{p} tiene dirección base \(A\), la expresión \texttt{p + k} (y su equivalente por definición \texttt{\&(p[k])}) evalúa a la dirección \(A + k \cdot \text{sizeof}(T)\). El estándar ISO C permite formar y comparar un puntero únicamente dentro del mismo array asignado o exactamente una posición más allá del último elemento (\texttt{p + N}, denominado *one-past-the-end*), aunque desreferenciar \texttt{p + N} es comportamiento indefinido. Las cadenas en C son secuencias de bytes terminadas por el carácter nulo \texttt{'\textbackslash{}0'} (\texttt{0x00}); calcular \texttt{strlen(s)} cuesta \(O(n)\) porque debe recorrer los bytes hasta hallar el centinela, y usar funciones no acotadas permite desbordar el búfer si no se pasa explícitamente la capacidad máxima del destino.
% studium:end-record eyJpZCI6ICJQQVItMDAwOCIsICJzZWN0aW9uIjogImNhcDAyLWMtcHVudGVyb3MtbWVtb3JpYS1kaW5hbWljYS1lc3RydWN0dXJhcy1hbGlhc2luZyJ9

\section{Bloque E: Punteros a función, calificador restrict, strict aliasing y asignadores de región (Arena)}
% studium:record eyJpZCI6ICJQQVItMDAwOSIsICJzZWN0aW9uIjogImNhcDAyLWMtcHVudGVyb3MtbWVtb3JpYS1kaW5hbWljYS1lc3RydWN0dXJhcy1hbGlhc2luZyJ9

La regla de *alias estricto* (*strict aliasing*, ISO C17 §6.5 p7) estipula que un objeto en memoria solo puede ser accedido a través de un lvalue de tipo compatible con su tipo efectivo o a través de un puntero a carácter (\texttt{unsigned char *} o \texttt{char *}). Por tanto, convertir un \texttt{uint8\_t *buf} arbitrario mediante \texttt{*(uint32\_t *)buf} viola tanto el alias estricto como potencialmente el alineamiento de hardware; la forma conforme al estándar \(-\)que GCC y Clang compilan a una única instrucción máquina \texttt{mov} sin llamada real a función\(-\) consiste en usar \texttt{memcpy(\&val, buf, sizeof(val))}. Asimismo, el calificador \texttt{restrict} en \texttt{void saxpy(size\_t n, float a, const float *restrict x, float *restrict y)} promete al compilador que las regiones apuntadas por \texttt{x} e \texttt{y} no se solapan, habilitando la vectorización SIMD automática:
\begin{lstlisting}[style=studium,language={C}]
#include <stddef.h>
#include <stdint.h>

void vector_scale_add(size_t n, float alpha, const float *restrict x, float *restrict y) {
    for (size_t i = 0; i < n; ++i) {
        y[i] += alpha * x[i];
    }
}
\end{lstlisting}
En sistemas de alto rendimiento (compiladores, servidores HTTP, videojuegos), en lugar de invocar miles de veces \texttt{malloc} y \texttt{free}, se reserva un bloque contiguo y se sirve memoria alineada incrementando un desplazamiento (\texttt{Arena Allocator}), liberando toda la región en \(O(1)\) al finalizar la petición o fase.
% studium:end-record eyJpZCI6ICJQQVItMDAwOSIsICJzZWN0aW9uIjogImNhcDAyLWMtcHVudGVyb3MtbWVtb3JpYS1kaW5hbWljYS1lc3RydWN0dXJhcy1hbGlhc2luZyJ9

\begin{tcolorbox}[title={Problema resuelto}, studiumworked, breakable]
% studium:record eyJpZCI6ICJQUkItMDAwNyIsICJzZWN0aW9uIjogImNhcDAyLWMtcHVudGVyb3MtbWVtb3JpYS1kaW5hbWljYS1lc3RydWN0dXJhcy1hbGlhc2luZyJ9

\noindent\textbf{Enunciado}
En el asignador por región \texttt{projects/c-arena-allocator}, el desplazamiento actual del búfer es \texttt{offset = 19} bytes y se solicita reservar una estructura de \(40\text{ B}\) con requisito de alineamiento \texttt{align = 8} bytes. El asignador avanza primero \texttt{offset} al siguiente múltiplo de 8 (sumando \(5\text{ B}\) de relleno hasta \(24\)) y luego suma los \(40\text{ B}\) del objeto. Calcula el nuevo valor exacto de \texttt{offset} tras la asignación.

\noindent\textbf{Datos}
\begin{itemize}
\item $offset_{0} = 19\,\mathrm{B}$ — Desplazamiento actual antes de la reserva
\item $pad = 5\,\mathrm{B}$ — Relleno para alcanzar el múltiplo de 8 (24 B)
\item $size = 40\,\mathrm{B}$ — Tamaño del bloque solicitado
\end{itemize}

\noindent\textbf{Incógnita}
Magnitud objetivo solicitada en el enunciado

\noindent\textbf{Modelo físico}

\noindent\textit{Avance de puntero con alineamiento en Arena Allocator}

\[
\text{offset}_1 = (19 + 5) + 40
\]

\noindent\textbf{Hipótesis}
\begin{itemize}
\item Evaluación determinista bajo el modelo formal del capítulo
\end{itemize}

\noindent\textbf{Desarrollo}
\begin{enumerate}
\item
Alineamos 19 + 5 = 24 y sumamos los 40 bytes del bloque

\[
\text{offset}_1 = (19 + 5) + 40
\]

\[
\left(\left(19+5\right)+40\right)
\]
\end{enumerate}

\noindent\textbf{Resultado}

\[
Resultado = 64. El desplazamiento final de la arena queda en 64 bytes.
\]

\noindent\textbf{Interpretación}
El desplazamiento final de la arena queda en 64 bytes.

\noindent\textbf{Comprobación de los pasos}
\begin{itemize}
\item Paso 1: $64$ — reproducido
\end{itemize}

\noindent\textbf{Comprobación de la respuesta} la aritmética del último paso reproduce la respuesta enunciada a la precisión con que está escrita. Es una reproducción aritmética; no valida el modelo físico ni sus hipótesis.
% studium:end-record eyJpZCI6ICJQUkItMDAwNyIsICJzZWN0aW9uIjogImNhcDAyLWMtcHVudGVyb3MtbWVtb3JpYS1kaW5hbWljYS1lc3RydWN0dXJhcy1hbGlhc2luZyJ9
\end{tcolorbox}

\begin{tcolorbox}[title={Ejercicios}, studiumworked, breakable]

\noindent\textbf{1 [PRB-0009] (Básico, Numérico)}
Si un puntero \texttt{uint64\_t *p} apunta a la dirección decimal 4096, ¿a qué dirección decimal apunta \texttt{p + 5} sabiendo que \texttt{sizeof(uint64\_t) = 8}?

\noindent\textbf{2 [PRB-0010] (Básico, Numérico)}
Una cadena en C contiene 31 caracteres visibles más su byte terminador \texttt{'\textbackslash{}0'}. ¿Cuántos bytes mínimos debe reservar \texttt{malloc}?

\noindent\textbf{3 [PRB-0011] (Intermedio, Desarrollo por pasos)}
En una arena con capacidad total de 4096 bytes, se reservan 12 nodos de 64 bytes ya alineados. ¿Cuántos bytes libres quedan en la arena?

\noindent\textbf{4 [PRB-0012] (Intermedio, Diseño)}
Si un paquete binario de telemetría tiene 8 bytes de cabecera y transporta 16 muestras \texttt{int16\_t} (2 bytes cada una), ¿cuántos bytes mide el paquete completo?
\end{tcolorbox}

\chapter{Monografía de C (III \(-\) Bloques F, G y H): Concurrencia C11/POSIX, Sanitizers, Ecosistema y Proyectos}
\label{chap:cap03-c-experto-atomicos-posix-sanitizers-ecosistema-proyectos}
% studium:record eyJpZCI6ICJQQVItMDAxMyIsICJzZWN0aW9uIjogImNhcDAzLWMtZXhwZXJ0by1hdG9taWNvcy1wb3NpeC1zYW5pdGl6ZXJzLWVjb3Npc3RlbWEtcHJveWVjdG9zIn0=

\noindent\textit{Este capítulo cierra la monografía de C abordando el nivel experto (Bloque F), el ecosistema de construcción y auditoría (Bloque G) y los proyectos prácticos (Bloque H), incluyendo la verificación real del proyecto \texttt{projects/c-arena-allocator} compilado con \texttt{clang -std=c17 -fsanitize=address,undefined}.}
% studium:end-record eyJpZCI6ICJQQVItMDAxMyIsICJzZWN0aW9uIjogImNhcDAzLWMtZXhwZXJ0by1hdG9taWNvcy1wb3NpeC1zYW5pdGl6ZXJzLWVjb3Npc3RlbWEtcHJveWVjdG9zIn0=

\begin{tcolorbox}[title={Consejo}, studiumtip, breakable]
% studium:record eyJpZCI6ICJQQVItMDAxNyIsICJzZWN0aW9uIjogImNhcDAzLWMtZXhwZXJ0by1hdG9taWNvcy1wb3NpeC1zYW5pdGl6ZXJzLWVjb3Npc3RlbWEtcHJveWVjdG9zIn0=

Nunca utilices \texttt{volatile} como sustituto de \texttt{<stdatomic.h>} o \texttt{pthread\_mutex\_t}: una carrera de datos sobre una variable \texttt{volatile} no atómica sigue siendo comportamiento indefinido según ISO C11/C17.
% studium:end-record eyJpZCI6ICJQQVItMDAxNyIsICJzZWN0aW9uIjogImNhcDAzLWMtZXhwZXJ0by1hdG9taWNvcy1wb3NpeC1zYW5pdGl6ZXJzLWVjb3Npc3RlbWEtcHJveWVjdG9zIn0=
\end{tcolorbox}

\begin{tcolorbox}[title={Definición}, studiumdef, breakable]
% studium:record eyJpZCI6ICJQQVItMDAxNiIsICJzZWN0aW9uIjogImNhcDAzLWMtZXhwZXJ0by1hdG9taWNvcy1wb3NpeC1zYW5pdGl6ZXJzLWVjb3Npc3RlbWEtcHJveWVjdG9zIn0=

Memoria sombra en AddressSanitizer (ASan Shadow Memory): técnica de instrumentación de compilador donde cada bloque alineado de 8 bytes de la memoria de aplicación se proyecta mediante \texttt{(addr >> 3) + Offset} sobre 1 byte de memoria sombra que codifica si dichos 8 bytes son direccionables o pertenecen a una zona roja (*redzone*) envenenada.
% studium:end-record eyJpZCI6ICJQQVItMDAxNiIsICJzZWN0aW9uIjogImNhcDAzLWMtZXhwZXJ0by1hdG9taWNvcy1wb3NpeC1zYW5pdGl6ZXJzLWVjb3Npc3RlbWEtcHJveWVjdG9zIn0=
\end{tcolorbox}

\begin{tcolorbox}[title={Autoficha}, studiumcheck, breakable]
% studium:record eyJpZCI6ICJQQVItMDAxOCIsICJzZWN0aW9uIjogImNhcDAzLWMtZXhwZXJ0by1hdG9taWNvcy1wb3NpeC1zYW5pdGl6ZXJzLWVjb3Npc3RlbWEtcHJveWVjdG9zIn0=

¿Por qué el patrón \texttt{goto cleanup;} hacia adelante al final de una función en C mejora la seguridad de liberación de recursos frente a tener múltiples sentencias \texttt{return} dispersas que olvidan liberar descriptores previos?
% studium:end-record eyJpZCI6ICJQQVItMDAxOCIsICJzZWN0aW9uIjogImNhcDAzLWMtZXhwZXJ0by1hdG9taWNvcy1wb3NpeC1zYW5pdGl6ZXJzLWVjb3Npc3RlbWEtcHJveWVjdG9zIn0=
\end{tcolorbox}

\section{Bloque F: Concurrencia C11 (<stdatomic.h>), pthreads, llamadas POSIX y programación defensiva}
% studium:record eyJpZCI6ICJQQVItMDAxNCIsICJzZWN0aW9uIjogImNhcDAzLWMtZXhwZXJ0by1hdG9taWNvcy1wb3NpeC1zYW5pdGl6ZXJzLWVjb3Npc3RlbWEtcHJveWVjdG9zIn0=

Un error histórico frecuente en C consiste en emplear el calificador \texttt{volatile} creyendo que proporciona sincronización entre hilos. En el estándar ISO C, \texttt{volatile} únicamente instruye al compilador para no omitir ni reordenar lecturas/escrituras respecto a registros de hardware mapeados en memoria (MMIO) o manejadores de señales (\texttt{sig\_atomic\_t}), pero **no** garantiza atomicidad ni emite barreras de coherencia de caché entre núcleos de CPU. Para concurrencia multihilo real, C11 introdujo \texttt{<stdatomic.h>} (\texttt{\_Atomic uint64\_t}, \texttt{atomic\_fetch\_add\_explicit}, \texttt{memory\_order\_acquire}, \texttt{memory\_order\_release}). En sistemas UNIX/POSIX (IEEE Std 1003.1), la gestión de recursos que pueden fallar en cascada dentro de una misma función se estructura idiomáticamente mediante el patrón de limpieza centralizada con etiquetas \texttt{goto cleanup:}:
\begin{lstlisting}[style=studium,language={C},escapeinside={(*@}{@*)}]
int process_file_buffer(size_t bytes) {
    int rc = -1;
    uint8_t *buf_a = malloc(bytes);
    if (!buf_a) goto fail_a;
    uint8_t *buf_b = malloc(bytes);
    if (!buf_b) goto fail_b;
    rc = 0; // operaci(*@{\'o}@*)n exitosa
    free(buf_b);
fail_b:
    free(buf_a);
fail_a:
    return rc;
}
\end{lstlisting}
Este patrón asegura un único camino determinista de liberación en orden inverso de adquisición sin duplicar llamadas a \texttt{free}.
% studium:end-record eyJpZCI6ICJQQVItMDAxNCIsICJzZWN0aW9uIjogImNhcDAzLWMtZXhwZXJ0by1hdG9taWNvcy1wb3NpeC1zYW5pdGl6ZXJzLWVjb3Npc3RlbWEtcHJveWVjdG9zIn0=

\section{Bloques G y H: Ecosistema (CMake, Meson, pkg-config, Sanitizers) y los 4 proyectos en C}
% studium:record eyJpZCI6ICJQQVItMDAxNSIsICJzZWN0aW9uIjogImNhcDAzLWMtZXhwZXJ0by1hdG9taWNvcy1wb3NpeC1zYW5pdGl6ZXJzLWVjb3Npc3RlbWEtcHJveWVjdG9zIn0=

El ecosistema moderno de C se apoya en generadores de sistemas de construcción como \texttt{CMake} y \texttt{Meson} sobre el ejecutor \texttt{Ninja}, el sistema de metadatos de bibliotecas \texttt{pkg-config}, pruebas unitarias ligeras y, de forma obligatoria en integración continua, la tríada de sanitizadores de compilador: \texttt{AddressSanitizer} (detecta desbordamientos en pila/montículo/globales y *use-after-free* con una sobrecarga media inferior a \(2\times\)), \texttt{UndefinedBehaviorSanitizer} (detecta desbordamientos con signo, desplazamientos ilegales y punteros nulos o desalineados) y \texttt{ThreadSanitizer} (detecta carreras de datos). Los cuatro proyectos del bloque en C son: (1) una utilidad CLI compatible con POSIX (\texttt{wc}/\texttt{sha256}) usando E/S por bloques; (2) el asignador por arena y analizador de paquetes binarios con checksum \texttt{projects/c-arena-allocator}; (3) una tabla hash abierta genérica con sondeo lineal y factor de carga controlado; y (4) un servidor TCP multicliente no bloqueante basado en multiplexación de descriptores.
% studium:end-record eyJpZCI6ICJQQVItMDAxNSIsICJzZWN0aW9uIjogImNhcDAzLWMtZXhwZXJ0by1hdG9taWNvcy1wb3NpeC1zYW5pdGl6ZXJzLWVjb3Npc3RlbWEtcHJveWVjdG9zIn0=

\begin{tcolorbox}[title={Problema resuelto}, studiumworked, breakable]
% studium:record eyJpZCI6ICJQUkItMDAxMyIsICJzZWN0aW9uIjogImNhcDAzLWMtZXhwZXJ0by1hdG9taWNvcy1wb3NpeC1zYW5pdGl6ZXJzLWVjb3Npc3RlbWEtcHJveWVjdG9zIn0=

\noindent\textbf{Enunciado}
Cuando un programa en C se compila con \texttt{-fsanitize=address} (AddressSanitizer), cada \(8\text{ B}\) de memoria de aplicación requieren \(1\text{ B}\) de memoria sombra (*shadow memory*), más zonas rojas (*redzones*) de \(32\text{ B}\) alrededor de cada asignación. Si una aplicación reserva \(N = 1000\) bloques de \(64\text{ B}\) cada uno (con \(32\text{ B}\) de zona roja por bloque), calcula cuántos bytes suman los bloques más sus zonas rojas más la memoria sombra correspondiente (\$1/8\$ del total).

\noindent\textbf{Datos}
\begin{itemize}
\item $N = 1000$ — Número de bloques asignados
\item $S_{bloque} = 64\,\mathrm{B}$ — Tamaño útil de cada bloque
\item $S_{redzone} = 32\,\mathrm{B}$ — Zona roja envenenada por bloque
\end{itemize}

\noindent\textbf{Incógnita}
Magnitud objetivo solicitada en el enunciado

\noindent\textbf{Modelo físico}

\noindent\textit{Consumo de memoria bajo instrumentación AddressSanitizer (ASan)}

\[
M = 1000 \cdot (64 + 32) \cdot \frac{9}{8}
\]

\noindent\textbf{Hipótesis}
\begin{itemize}
\item Evaluación determinista bajo el modelo formal del capítulo
\end{itemize}

\noindent\textbf{Desarrollo}
\begin{enumerate}
\item
Sumamos bloque + redzone (96000 B) y multiplicamos por 9/8 para incluir el byte sombra por cada 8 B

\[
M = 1000 \cdot (64 + 32) \cdot \frac{9}{8}
\]

\[
\frac{\left(\left(1000\cdot \left(64+32\right)\right)\cdot 9\right)}{8}
\]
\end{enumerate}

\noindent\textbf{Resultado}

\[
Resultado = 108000. El conjunto instrumentado ocupa 108 000 bytes (96 000 B de bloques con redzones y 12 000 B de memoria sombra).
\]

\noindent\textbf{Interpretación}
El conjunto instrumentado ocupa 108 000 bytes (96 000 B de bloques con redzones y 12 000 B de memoria sombra).

\noindent\textbf{Comprobación de los pasos}
\begin{itemize}
\item Paso 1: $108000$ — reproducido
\end{itemize}

\noindent\textbf{Comprobación de la respuesta} la aritmética del último paso reproduce la respuesta enunciada a la precisión con que está escrita. Es una reproducción aritmética; no valida el modelo físico ni sus hipótesis.
% studium:end-record eyJpZCI6ICJQUkItMDAxMyIsICJzZWN0aW9uIjogImNhcDAzLWMtZXhwZXJ0by1hdG9taWNvcy1wb3NpeC1zYW5pdGl6ZXJzLWVjb3Npc3RlbWEtcHJveWVjdG9zIn0=
\end{tcolorbox}

\begin{tcolorbox}[title={Ejercicios}, studiumworked, breakable]

\noindent\textbf{1 [PRB-0015] (Básico, Numérico)}
En AddressSanitizer, ¿cuántos bytes de memoria sombra corresponden a 64 000 bytes de memoria de aplicación (\$64000 / 8\$)?

\noindent\textbf{2 [PRB-0016] (Básico, Numérico)}
En \texttt{projects/c-arena-allocator}, la cabecera de telemetría suma los campos \texttt{magic} (10), \texttt{version} (2), \texttt{flags} (3) y \texttt{payload\_len} (25). ¿Cuál es el valor del checksum esperado?

\noindent\textbf{3 [PRB-0017] (Intermedio, Desarrollo por pasos)}
Una función en C adquiere secuencialmente 5 recursos y falla al intentar adquirir el cuarto. ¿Cuántos recursos ya adquiridos debe liberar la cascada \texttt{goto cleanup} antes de retornar?

\noindent\textbf{4 [PRB-0018] (Intermedio, Diseño)}
Un anillo de descriptores E/S tiene capacidad para 256 entradas de 16 bytes cada una. ¿Cuántos bytes mide el anillo completo mapeado en memoria?
\end{tcolorbox}

\part{Monografía Autónoma de C++ (ISO/IEC 14882:2020 / C++23)}

\chapter{Monografía de C++ (I \(-\) Bloques A\(-\)D): De C++98 a C++20/23, RAII, Semántica de Movimiento y Punteros Inteligentes}
\label{chap:cap04-cpp-fundamentos-raii-move-semantics-smart-pointers}
% studium:record eyJpZCI6ICJQQVItMDAxOSIsICJzZWN0aW9uIjogImNhcDA0LWNwcC1mdW5kYW1lbnRvcy1yYWlpLW1vdmUtc2VtYW50aWNzLXNtYXJ0LXBvaW50ZXJzIn0=

\noindent\textit{Este capítulo inicia la monografía autónoma de C++ moderno (ISO C++20 / C++23). Lejos de tratar a C++ como «C con clases y \texttt{new}/\texttt{delete} manuales», enseñamos desde el primer día el paradigma moderno: propiedad automática determinista mediante RAII, vistas no propietarias (\texttt{std::string\_view}, \texttt{std::span}), evaluación en tiempo de compilación (\texttt{constexpr}/\texttt{consteval}) y semántica de movimiento.}
% studium:end-record eyJpZCI6ICJQQVItMDAxOSIsICJzZWN0aW9uIjogImNhcDA0LWNwcC1mdW5kYW1lbnRvcy1yYWlpLW1vdmUtc2VtYW50aWNzLXNtYXJ0LXBvaW50ZXJzIn0=

\begin{tcolorbox}[title={Consejo}, studiumtip, breakable]
% studium:record eyJpZCI6ICJQQVItMDAyMyIsICJzZWN0aW9uIjogImNhcDA0LWNwcC1mdW5kYW1lbnRvcy1yYWlpLW1vdmUtc2VtYW50aWNzLXNtYXJ0LXBvaW50ZXJzIn0=

En C++20/C++23, no escribas jamás \texttt{new} ni \texttt{delete} desnudos en código de aplicación: utiliza \texttt{std::make\_unique<T>(...)} para propiedad única (cuyo tamaño y coste en ensamblador son idénticos a un puntero crudo \texttt{T*}) y pasa vistas \texttt{std::span<const T>} o \texttt{std::string\_view} cuando una función no necesite adquirir propiedad.
% studium:end-record eyJpZCI6ICJQQVItMDAyMyIsICJzZWN0aW9uIjogImNhcDA0LWNwcC1mdW5kYW1lbnRvcy1yYWlpLW1vdmUtc2VtYW50aWNzLXNtYXJ0LXBvaW50ZXJzIn0=
\end{tcolorbox}

\begin{tcolorbox}[title={Definición}, studiumdef, breakable]
% studium:record eyJpZCI6ICJQQVItMDAyMiIsICJzZWN0aW9uIjogImNhcDA0LWNwcC1mdW5kYW1lbnRvcy1yYWlpLW1vdmUtc2VtYW50aWNzLXNtYXJ0LXBvaW50ZXJzIn0=

Garantía fuerte frente a excepciones (Strong Exception Safety Guarantee): propiedad de una operación en C++ según la cual, si la operación termina lanzando una excepción, el estado observable del programa permanece idéntico al que tenía antes de iniciar la operación (semántica commit-or-rollback).
% studium:end-record eyJpZCI6ICJQQVItMDAyMiIsICJzZWN0aW9uIjogImNhcDA0LWNwcC1mdW5kYW1lbnRvcy1yYWlpLW1vdmUtc2VtYW50aWNzLXNtYXJ0LXBvaW50ZXJzIn0=
\end{tcolorbox}

\begin{tcolorbox}[title={Autoficha}, studiumcheck, breakable]
% studium:record eyJpZCI6ICJQQVItMDAyNCIsICJzZWN0aW9uIjogImNhcDA0LWNwcC1mdW5kYW1lbnRvcy1yYWlpLW1vdmUtc2VtYW50aWNzLXNtYXJ0LXBvaW50ZXJzIn0=

¿Por qué \texttt{std::vector<T>::push\_back} se ve obligado a copiar en lugar de mover los elementos existentes al duplicar su capacidad si el constructor de movimiento \texttt{T(T\&\&)} no está declarado como \texttt{noexcept}?
% studium:end-record eyJpZCI6ICJQQVItMDAyNCIsICJzZWN0aW9uIjogImNhcDA0LWNwcC1mdW5kYW1lbnRvcy1yYWlpLW1vdmUtc2VtYW50aWNzLXNtYXJ0LXBvaW50ZXJzIn0=
\end{tcolorbox}

\section{Bloques A\(-\)C: Evolución C++11/14/17/20/23, toolchain, referencias, sobrecarga, constexpr y RAII}
% studium:record eyJpZCI6ICJQQVItMDAyMCIsICJzZWN0aW9uIjogImNhcDA0LWNwcC1mdW5kYW1lbnRvcy1yYWlpLW1vdmUtc2VtYW50aWNzLXNtYXJ0LXBvaW50ZXJzIn0=

Diseñado por Bjarne Stroustrup desde 1979 y estandarizado por ISO en 1998, C++ experimentó una profunda modernización con C++11 (semántica de movimiento, lambdas, \texttt{auto}, \texttt{constexpr}, \texttt{std::unique\_ptr}, modelo de memoria multihilo), C++14, C++17 (\texttt{std::optional}, \texttt{std::variant}, \texttt{std::string\_view}, \texttt{if constexpr}) y C++20/C++23 (\texttt{concepts}, \texttt{std::ranges}, corrutinas, módulos, \texttt{std::span}, \texttt{std::expected}). El pilar arquitectónico de C++ es **RAII** (*Resource Acquisition Is Initialization*): todo recurso (memoria del montículo, socket, fichero, cerrojo \texttt{std::scoped\_lock}) se vincula al ciclo de vida de un objeto alojado en la pila cuyo constructor establece la invariante y cuyo destructor \texttt{\textasciitilde{}T()} libera infaliblemente el recurso al salir del ámbito, incluso si se desenrolla la pila (*stack unwinding*) por el lanzamiento de una excepción.
% studium:end-record eyJpZCI6ICJQQVItMDAyMCIsICJzZWN0aW9uIjogImNhcDA0LWNwcC1mdW5kYW1lbnRvcy1yYWlpLW1vdmUtc2VtYW50aWNzLXNtYXJ0LXBvaW50ZXJzIn0=

\section{Bloque D: Categorías de valores (lvalue/rvalue), Regla de Cero/Cinco, std::unique\_ptr y std::span}
% studium:record eyJpZCI6ICJQQVItMDAyMSIsICJzZWN0aW9uIjogImNhcDA0LWNwcC1mdW5kYW1lbnRvcy1yYWlpLW1vdmUtc2VtYW50aWNzLXNtYXJ0LXBvaW50ZXJzIn0=

C++11 clasificó toda expresión según su identidad y movibilidad en \texttt{lvalue} (tiene identidad persistente), \texttt{xvalue} (*expiring value*, tiene identidad pero sus recursos pueden ser robados) y \texttt{prvalue} (temporal puro). La **Regla de Cero** establece que las clases de lógica de negocio no deben escribir destructores, constructores de copia ni de movimiento manuales, sino componer miembros RAII (\texttt{std::vector}, \texttt{std::string}, \texttt{std::unique\_ptr}). Si una clase gestiona directamente un manejador crudo del sistema, aplica la **Regla de Cinco**: definir explícitamente destructor, constructor de copia, asignación por copia, constructor de movimiento (\texttt{noexcept}) y asignación por movimiento (\texttt{noexcept}):
\begin{lstlisting}[style=studium,language={C++}]
#include <memory>
#include <span>
#include <numeric>

// std::span<const double> observa cualquier array, std::vector o std::array sin copiar:
constexpr double sum_window(std::span<const double> samples) noexcept {
    double acc = 0.0;
    for (double v : samples) { acc += v; }
    return acc;
}
\end{lstlisting}
Es fundamental marcar los constructores de movimiento como \texttt{noexcept}: \texttt{std::vector} solo mueve sus elementos durante una realojación si el constructor de movimiento promete \texttt{noexcept}; de lo contrario, cae de vuelta a copiar cada elemento para preservar la garantía fuerte frente a excepciones (*strong exception safety guarantee*).
% studium:end-record eyJpZCI6ICJQQVItMDAyMSIsICJzZWN0aW9uIjogImNhcDA0LWNwcC1mdW5kYW1lbnRvcy1yYWlpLW1vdmUtc2VtYW50aWNzLXNtYXJ0LXBvaW50ZXJzIn0=

\begin{tcolorbox}[title={Problema resuelto}, studiumworked, breakable]
% studium:record eyJpZCI6ICJQUkItMDAxOSIsICJzZWN0aW9uIjogImNhcDA0LWNwcC1mdW5kYW1lbnRvcy1yYWlpLW1vdmUtc2VtYW50aWNzLXNtYXJ0LXBvaW50ZXJzIn0=

\noindent\textbf{Enunciado}
En una implementación estándar de C++ de 64 bits donde \texttt{sizeof(std::string) = 24} bytes y la optimización de cadena corta (*Small String Optimization*, SSO) almacena en línea cadenas de hasta \(22\text{ B}\) sin alojar en el montículo, un programa crea un \texttt{std::vector<std::string>} con \(N = 5000\) identificadores cortos de \(12\text{ B}\) cada uno. Calcula cuántos bytes totales ocupan los 5000 objetos \texttt{std::string} dentro del búfer contiguo del vector (sabiendo que realizan \(0\) alojamientos adicionales fuera del vector).

\noindent\textbf{Datos}
\begin{itemize}
\item $N = 5000$ — Número de cadenas cortas de 12 B (<= 22 B SSO)
\item $S_{str} = 24\,\mathrm{B}$ — Tamaño de cada objeto std::string en 64 bits
\end{itemize}

\noindent\textbf{Incógnita}
Magnitud objetivo solicitada en el enunciado

\noindent\textbf{Modelo físico}

\noindent\textit{Huella de memoria con Small String Optimization (SSO) en C++}

\[
M = 5000 \cdot 24
\]

\noindent\textbf{Hipótesis}
\begin{itemize}
\item Evaluación determinista bajo el modelo formal del capítulo
\end{itemize}

\noindent\textbf{Desarrollo}
\begin{enumerate}
\item
Multiplicamos las 5000 instancias por 24 B (sin alojamientos externos gracias a SSO)

\[
M = 5000 \cdot 24
\]

\[
\left(5000\cdot 24\right)
\]
\end{enumerate}

\noindent\textbf{Resultado}

\[
Resultado = 120000. Los 5000 identificadores se almacenan contiguamente en 120 000 bytes con cero llamadas individuales al montículo por cadena.
\]

\noindent\textbf{Interpretación}
Los 5000 identificadores se almacenan contiguamente en 120 000 bytes con cero llamadas individuales al montículo por cadena.

\noindent\textbf{Comprobación de los pasos}
\begin{itemize}
\item Paso 1: $120000$ — reproducido
\end{itemize}

\noindent\textbf{Comprobación de la respuesta} la aritmética del último paso reproduce la respuesta enunciada a la precisión con que está escrita. Es una reproducción aritmética; no valida el modelo físico ni sus hipótesis.
% studium:end-record eyJpZCI6ICJQUkItMDAxOSIsICJzZWN0aW9uIjogImNhcDA0LWNwcC1mdW5kYW1lbnRvcy1yYWlpLW1vdmUtc2VtYW50aWNzLXNtYXJ0LXBvaW50ZXJzIn0=
\end{tcolorbox}

\begin{tcolorbox}[title={Ejercicios}, studiumworked, breakable]

\noindent\textbf{1 [PRB-0021] (Básico, Numérico)}
¿Cuántos bytes ocupa \texttt{std::unique\_ptr<double>} con deleción por defecto en una arquitectura de 64 bits?

\noindent\textbf{2 [PRB-0022] (Básico, Numérico)}
Un \texttt{std::span<const int>} en C++20 con extensión dinámica almacena un puntero (8 B) y una longitud \texttt{size\_t} (8 B). ¿Cuántos bytes ocupa en la pila?

\noindent\textbf{3 [PRB-0023] (Intermedio, Desarrollo por pasos)}
Un \texttt{std::vector<int>} duplica su capacidad geométricamente desde 16 hasta 1024 elementos (\texttt{sizeof(int) = 4}). ¿Cuántos bytes reserva el búfer final de 1024 enteros?

\noindent\textbf{4 [PRB-0024] (Intermedio, Diseño)}
Un bloque de control de \texttt{std::make\_shared<T>} combina en una única reserva de montículo los dos contadores atómicos (16 B) más el objeto \texttt{T} de 48 B. ¿Cuántos bytes mide esa única reserva?
\end{tcolorbox}

\chapter{Monografía de C++ (II \(-\) Bloques E\(-\)H): Plantillas, Concepts C++20, std::ranges, Concurrencia, Corrutinas y Proyectos}
\label{chap:cap05-cpp-templates-concepts-ranges-concurrencia-proyectos}
% studium:record eyJpZCI6ICJQQVItMDAyNSIsICJzZWN0aW9uIjogImNhcDA1LWNwcC10ZW1wbGF0ZXMtY29uY2VwdHMtcmFuZ2VzLWNvbmN1cnJlbmNpYS1wcm95ZWN0b3MifQ==

\noindent\textit{Este capítulo completa la monografía de C++ desarrollando los Bloques E, F, G y H: desde la programación genérica con \texttt{concept}s y \texttt{std::ranges} hasta la concurrencia cooperativa con \texttt{std::jthread} y la verificación del planificador de tareas \texttt{projects/cpp-task-scheduler}.}
% studium:end-record eyJpZCI6ICJQQVItMDAyNSIsICJzZWN0aW9uIjogImNhcDA1LWNwcC10ZW1wbGF0ZXMtY29uY2VwdHMtcmFuZ2VzLWNvbmN1cnJlbmNpYS1wcm95ZWN0b3MifQ==

\begin{tcolorbox}[title={Consejo}, studiumtip, breakable]
% studium:record eyJpZCI6ICJQQVItMDAyOSIsICJzZWN0aW9uIjogImNhcDA1LWNwcC10ZW1wbGF0ZXMtY29uY2VwdHMtcmFuZ2VzLWNvbmN1cnJlbmNpYS1wcm95ZWN0b3MifQ==

Sustituye \texttt{std::thread} por \texttt{std::jthread} en todo código C++20 nuevo: si una función lanza una excepción antes de llamar a \texttt{.join()} sobre un \texttt{std::thread} convencional, su destructor invoca \texttt{std::terminate()} abortando el proceso.
% studium:end-record eyJpZCI6ICJQQVItMDAyOSIsICJzZWN0aW9uIjogImNhcDA1LWNwcC10ZW1wbGF0ZXMtY29uY2VwdHMtcmFuZ2VzLWNvbmN1cnJlbmNpYS1wcm95ZWN0b3MifQ==
\end{tcolorbox}

\begin{tcolorbox}[title={Definición}, studiumdef, breakable]
% studium:record eyJpZCI6ICJQQVItMDAyOCIsICJzZWN0aW9uIjogImNhcDA1LWNwcC10ZW1wbGF0ZXMtY29uY2VwdHMtcmFuZ2VzLWNvbmN1cnJlbmNpYS1wcm95ZWN0b3MifQ==

Concepto en C++20 (\texttt{concept}): predicado nombrado evaluado en tiempo de compilación que restringe el conjunto de argumentos de plantilla aceptados por una función o clase genérica y participa en la resolución de sobrecarga por subsunción lógica.
% studium:end-record eyJpZCI6ICJQQVItMDAyOCIsICJzZWN0aW9uIjogImNhcDA1LWNwcC10ZW1wbGF0ZXMtY29uY2VwdHMtcmFuZ2VzLWNvbmN1cnJlbmNpYS1wcm95ZWN0b3MifQ==
\end{tcolorbox}

\begin{tcolorbox}[title={Autoficha}, studiumcheck, breakable]
% studium:record eyJpZCI6ICJQQVItMDAzMCIsICJzZWN0aW9uIjogImNhcDA1LWNwcC10ZW1wbGF0ZXMtY29uY2VwdHMtcmFuZ2VzLWNvbmN1cnJlbmNpYS1wcm95ZWN0b3MifQ==

¿Cómo decide el compilador de C++20 qué sobrecarga de plantilla es más específica cuando dos funciones restringidas por \texttt{concept}s distintos son válidas para un mismo tipo (regla de subsunción)?
% studium:end-record eyJpZCI6ICJQQVItMDAzMCIsICJzZWN0aW9uIjogImNhcDA1LWNwcC10ZW1wbGF0ZXMtY29uY2VwdHMtcmFuZ2VzLWNvbmN1cnJlbmNpYS1wcm95ZWN0b3MifQ==
\end{tcolorbox}

\section{Bloques E y F: Plantillas, C++20 Concepts, std::ranges, std::jthread y corrutinas}
% studium:record eyJpZCI6ICJQQVItMDAyNiIsICJzZWN0aW9uIjogImNhcDA1LWNwcC10ZW1wbGF0ZXMtY29uY2VwdHMtcmFuZ2VzLWNvbmN1cnJlbmNpYS1wcm95ZWN0b3MifQ==

Históricamente, las plantillas (*templates*) en C++ dependían de SFINAE (*Substitution Failure Is Not An Error*), produciendo diagnósticos de compilación crípticos cuando un tipo no cumplía una operación interna. C++20 resolvió este problema introduciendo **Concepts**: predicados en tiempo de compilación que especifican formalmente los requisitos sintácticos y semánticos de un parámetro de plantilla. En nuestro proyecto verificado \texttt{projects/cpp-task-scheduler}, el concepto \texttt{ExecutableTask} exige que cualquier tipo de tarea proporcione \texttt{id()}, \texttt{priority()} convertible a \texttt{int} y \texttt{execute()} convertible a \texttt{int64\_t}:
\begin{lstlisting}[style=studium,language={C++}]
#include <concepts>
#include <cstdint>

template <typename T>
concept ExecutableTask = requires(const T& t) {
    { t.id() } -> std::convertible_to<uint32_t>;
    { t.priority() } -> std::convertible_to<int>;
    { t.execute() } -> std::convertible_to<int64_t>;
};
\end{lstlisting}
Junto con los conceptos, la biblioteca \texttt{<ranges>} permite encadenar vistas perezosas mediante el operador tubería \texttt{|}, como \texttt{tasks | std::views::filter(...) | std::views::transform(...)}, y \texttt{<thread>} incorpora \texttt{std::jthread}, que se une automáticamente (\texttt{join()}) en su destructor e integra cancelación cooperativa mediante \texttt{std::stop\_token}.
% studium:end-record eyJpZCI6ICJQQVItMDAyNiIsICJzZWN0aW9uIjogImNhcDA1LWNwcC10ZW1wbGF0ZXMtY29uY2VwdHMtcmFuZ2VzLWNvbmN1cnJlbmNpYS1wcm95ZWN0b3MifQ==

\section{Bloques G y H: Ecosistema CMake/vcpkg/Conan/Clang-Tidy y los 4 proyectos (incluyendo cpp-task-scheduler)}
% studium:record eyJpZCI6ICJQQVItMDAyNyIsICJzZWN0aW9uIjogImNhcDA1LWNwcC10ZW1wbGF0ZXMtY29uY2VwdHMtcmFuZ2VzLWNvbmN1cnJlbmNpYS1wcm95ZWN0b3MifQ==

En el Bloque G (Ecosistema), los proyectos profesionales de C++ se estructuran mediante \texttt{CMake} moderno basado en objetivos (\texttt{target\_link\_libraries}, \texttt{target\_compile\_features(tgt PUBLIC cxx\_std\_20)}), gestores de paquetes \texttt{vcpkg} o \texttt{Conan 2}, pruebas con \texttt{GoogleTest} o \texttt{Catch2}, benchmarking con \texttt{Google Benchmark} y auditoría estática continua con \texttt{clang-tidy} (perfiles \texttt{cppcoreguidelines-*}, \texttt{modernize-*}, \texttt{bugprone-*}). En el Bloque H, además de una biblioteca de álgebra lineal estática \texttt{constexpr}, un motor de eventos con \texttt{std::variant} y un servidor de red asíncrono, el lector cuenta con el proyecto completo \texttt{projects/cpp-task-scheduler}, que implementa una cola de prioridad determinista \texttt{TaskScheduler<Task>} verificada con \texttt{clang++ -std=c++20 -fsanitize=address,undefined}.
% studium:end-record eyJpZCI6ICJQQVItMDAyNyIsICJzZWN0aW9uIjogImNhcDA1LWNwcC10ZW1wbGF0ZXMtY29uY2VwdHMtcmFuZ2VzLWNvbmN1cnJlbmNpYS1wcm95ZWN0b3MifQ==

\begin{tcolorbox}[title={Problema resuelto}, studiumworked, breakable]
% studium:record eyJpZCI6ICJQUkItMDAyNSIsICJzZWN0aW9uIjogImNhcDA1LWNwcC10ZW1wbGF0ZXMtY29uY2VwdHMtcmFuZ2VzLWNvbmN1cnJlbmNpYS1wcm95ZWN0b3MifQ==

\noindent\textbf{Enunciado}
En el proyecto \texttt{projects/cpp-task-scheduler}, se encolan cuatro tareas \texttt{(id, prioridad, base\_cost)}: \texttt{(1, 2, 10)}, \texttt{(2, 9, 10)}, \texttt{(3, 5, 20)} y \texttt{(4, 9, 15)}. Si se ejecuta \texttt{execute\_filtered(5)} (que filtra las tareas con \texttt{priority >= 5} y suma \texttt{base\_cost * priority} de cada una), calcula el resultado exacto devuelto.

\noindent\textbf{Datos}
\begin{itemize}
\item $T_{2} = prioridad 9, coste 10$ — Aporta 9 * 10 = 90
\item $T_{4} = prioridad 9, coste 15$ — Aporta 9 * 15 = 135
\item $T_{3} = prioridad 5, coste 20$ — Aporta 5 * 20 = 100
\end{itemize}

\noindent\textbf{Incógnita}
Magnitud objetivo solicitada en el enunciado

\noindent\textbf{Modelo físico}

\noindent\textit{Filtrado y reducción en cpp-task-scheduler con std::views}

\[
R = 9 \cdot 10 + 9 \cdot 15 + 5 \cdot 20
\]

\noindent\textbf{Hipótesis}
\begin{itemize}
\item Evaluación determinista bajo el modelo formal del capítulo
\end{itemize}

\noindent\textbf{Desarrollo}
\begin{enumerate}
\item
Descartamos la tarea 1 (prioridad 2 < 5) y sumamos los productos de las tres restantes

\[
R = 9 \cdot 10 + 9 \cdot 15 + 5 \cdot 20
\]

\[
\left(\left(\left(9\cdot 10\right)+\left(9\cdot 15\right)\right)+\left(5\cdot 20\right)\right)
\]
\end{enumerate}

\noindent\textbf{Resultado}

\[
Resultado = 325. El planificador devuelve exactamente 325, coincidiendo con el test ejecutable de `projects/cpp-task-scheduler`.
\]

\noindent\textbf{Interpretación}
El planificador devuelve exactamente 325, coincidiendo con el test ejecutable de \texttt{projects/cpp-task-scheduler}.

\noindent\textbf{Comprobación de los pasos}
\begin{itemize}
\item Paso 1: $325$ — reproducido
\end{itemize}

\noindent\textbf{Comprobación de la respuesta} la aritmética del último paso reproduce la respuesta enunciada a la precisión con que está escrita. Es una reproducción aritmética; no valida el modelo físico ni sus hipótesis.
% studium:end-record eyJpZCI6ICJQUkItMDAyNSIsICJzZWN0aW9uIjogImNhcDA1LWNwcC10ZW1wbGF0ZXMtY29uY2VwdHMtcmFuZ2VzLWNvbmN1cnJlbmNpYS1wcm95ZWN0b3MifQ==
\end{tcolorbox}

\begin{tcolorbox}[title={Ejercicios}, studiumworked, breakable]

\noindent\textbf{1 [PRB-0027] (Básico, Numérico)}
En \texttt{cpp-task-scheduler}, ¿qué valor devuelve \texttt{ComputeTask\{10, 7, 12\}.execute()} (\(12 \times 7\))?

\noindent\textbf{2 [PRB-0028] (Básico, Numérico)}
Una plantilla variádica \texttt{sum\_all(10, 20, 30, 40)} utiliza una expresión pliegue (*fold expression* \texttt{(args + ...)}). ¿Cuál es el resultado evaluado en \texttt{constexpr}?

\noindent\textbf{3 [PRB-0029] (Intermedio, Desarrollo por pasos)}
Un \texttt{std::variant<int, double, std::array<char, 32>>} almacena la alternativa mayor (32 B) más un índice alineado a 8 B. ¿Cuántos bytes ocupa la instancia?

\noindent\textbf{4 [PRB-0030] (Intermedio, Diseño)}
Si una vista \texttt{std::views::iota(1, 1001) | std::views::filter(...) | std::views::take(5)} detiene la evaluación al encontrar los primeros 5 pares (\$2, 4, 6, 8, 10\$), ¿cuántos enteros del rango \$[1, 1000]\$ llegaron a inspeccionarse?
\end{tcolorbox}

\part{Monografías Especializadas de Sistemas: Zig, Ensamblador, Ada y Fortran}

\chapter{Monografía de Zig (Bloques A\(-\)H): Asignadores explícitos, metaprogramación comptime, uniones de error y toolchain C}
\label{chap:cap06-monografia-zig-asignadores-comptime-interoperabilidad-c}
% studium:record eyJpZCI6ICJQQVItMDAzMSIsICJzZWN0aW9uIjogImNhcDA2LW1vbm9ncmFmaWEtemlnLWFzaWduYWRvcmVzLWNvbXB0aW1lLWludGVyb3BlcmFiaWxpZGFkLWMifQ==

\noindent\textit{Este capítulo presenta la monografía completa de **Zig** siguiendo los bloques \(A-H\). Creado por Andrew Kelley, Zig se concibe como un sucesor moderno y transparente de C donde ninguna expresión oculta llamadas a funciones ni reservas de memoria, los tipos son valores de primera clase en tiempo de compilación (\texttt{comptime}) y el propio binario \texttt{zig} incluye un compilador C/C++ de compilación cruzada basado en Clang/LLVM.}
% studium:end-record eyJpZCI6ICJQQVItMDAzMSIsICJzZWN0aW9uIjogImNhcDA2LW1vbm9ncmFmaWEtemlnLWFzaWduYWRvcmVzLWNvbXB0aW1lLWludGVyb3BlcmFiaWxpZGFkLWMifQ==

\begin{tcolorbox}[title={Consejo}, studiumtip, breakable]
% studium:record eyJpZCI6ICJQQVItMDAzNSIsICJzZWN0aW9uIjogImNhcDA2LW1vbm9ncmFmaWEtemlnLWFzaWduYWRvcmVzLWNvbXB0aW1lLWludGVyb3BlcmFiaWxpZGFkLWMifQ==

En todas tus pruebas unitarias de Zig (\texttt{test "..." \{ ... \}}), pasa siempre \texttt{std.testing.allocator} a las estructuras bajo prueba: el arnés de pruebas hará fallar automáticamente el test si queda un solo byte sin liberar o si ocurre doble liberación.
% studium:end-record eyJpZCI6ICJQQVItMDAzNSIsICJzZWN0aW9uIjogImNhcDA2LW1vbm9ncmFmaWEtemlnLWFzaWduYWRvcmVzLWNvbXB0aW1lLWludGVyb3BlcmFiaWxpZGFkLWMifQ==
\end{tcolorbox}

\begin{tcolorbox}[title={Definición}, studiumdef, breakable]
% studium:record eyJpZCI6ICJQQVItMDAzNCIsICJzZWN0aW9uIjogImNhcDA2LW1vbm9ncmFmaWEtemlnLWFzaWduYWRvcmVzLWNvbXB0aW1lLWludGVyb3BlcmFiaWxpZGFkLWMifQ==

Limpieza condicional de error (\texttt{errdefer} en Zig): sentencia registrada en la pila léxica del bloque actual cuya expresión asociada se ejecuta si y solo si la función abandona dicho ámbito devolviendo un valor de error de su unión de retorno \texttt{!T}.
% studium:end-record eyJpZCI6ICJQQVItMDAzNCIsICJzZWN0aW9uIjogImNhcDA2LW1vbm9ncmFmaWEtemlnLWFzaWduYWRvcmVzLWNvbXB0aW1lLWludGVyb3BlcmFiaWxpZGFkLWMifQ==
\end{tcolorbox}

\begin{tcolorbox}[title={Autoficha}, studiumcheck, breakable]
% studium:record eyJpZCI6ICJQQVItMDAzNiIsICJzZWN0aW9uIjogImNhcDA2LW1vbm9ncmFmaWEtemlnLWFzaWduYWRvcmVzLWNvbXB0aW1lLWludGVyb3BlcmFiaWxpZGFkLWMifQ==

¿Qué ventaja arquitectónica aporta que \texttt{std.ArrayList} en Zig reciba un \texttt{Allocator} explícito en lugar de invocar un asignador global fijo como \texttt{std::vector} por defecto en C++ o \texttt{Vec<T>} en Rust estable?
% studium:end-record eyJpZCI6ICJQQVItMDAzNiIsICJzZWN0aW9uIjogImNhcDA2LW1vbm9ncmFmaWEtemlnLWFzaWduYWRvcmVzLWNvbXB0aW1lLWludGVyb3BlcmFiaWxpZGFkLWMifQ==
\end{tcolorbox}

\section{Bloques A\(-\)D: Filosofía de Zig, toolchain integrada, tipos, rebanadas, uniones de error y defer/errdefer}
% studium:record eyJpZCI6ICJQQVItMDAzMiIsICJzZWN0aW9uIjogImNhcDA2LW1vbm9ncmFmaWEtemlnLWFzaWduYWRvcmVzLWNvbXB0aW1lLWludGVyb3BlcmFiaWxpZGFkLWMifQ==

Zig se rige por tres mandamientos de ingeniería de sistemas: (1) sin flujo de control oculto (no hay sobrecarga de operadores, ni constructores/destructores implícitos, ni excepciones); (2) sin asignaciones de memoria ocultas (toda función de la biblioteca estándar que necesita alojar en el montículo recibe explícitamente un parámetro \texttt{allocator: std.mem.Allocator}); y (3) sin preprocesador de macros textuales. Para gestionar recursos y errores con precisión quirúrgica, Zig combina tipos de unión de error (\texttt{AllocError![]u8}) con las palabras clave \texttt{try}, \texttt{catch}, \texttt{defer} (que se ejecuta al salir del bloque actual) y \texttt{errdefer} (que se ejecuta **únicamente** si el bloque actual retorna con un error):
\begin{lstlisting}[style=studium]
const std = @import("std");

pub fn allocateInitializedRing(allocator: std.mem.Allocator, n: usize) ![]u64 {
    const buf = try allocator.alloc(u64, n);
    errdefer allocator.free(buf);
    for (buf, 0..) |*slot, idx| {
        slot.* = @as(u64, @intCast(idx)) * 10;
    }
    return buf;
}
\end{lstlisting}
Si cualquier paso posterior dentro de la función fallara y retornara un error mediante \texttt{try}, \texttt{errdefer} libera automáticamente \texttt{buf} sin necesidad de etiquetas \texttt{goto}.
% studium:end-record eyJpZCI6ICJQQVItMDAzMiIsICJzZWN0aW9uIjogImNhcDA2LW1vbm9ncmFmaWEtemlnLWFzaWduYWRvcmVzLWNvbXB0aW1lLWludGVyb3BlcmFiaWxpZGFkLWMifQ==

\section{Bloques E\(-\)H: Asignadores explícitos (Arena, FixedBuffer, GPA), comptime, FFI con C y proyectos}
% studium:record eyJpZCI6ICJQQVItMDAzMyIsICJzZWN0aW9uIjogImNhcDA2LW1vbm9ncmFmaWEtemlnLWFzaWduYWRvcmVzLWNvbXB0aW1lLWludGVyb3BlcmFiaWxpZGFkLWMifQ==

En lugar de plantillas o macros, el polimorfismo paramétrico en Zig consiste en funciones ordinarias que aceptan argumentos \texttt{comptime T: type} y devuelven un \texttt{type} (\texttt{fn BoundedQueue(comptime T: type, comptime cap: usize) type}). Gracias a que los asignadores se pasan como valores de interfaz \texttt{std.mem.Allocator}, el mismo contenedor \texttt{std.ArrayList(T)} puede ejecutarse sobre el asignador detector de fugas \texttt{std.testing.allocator} durante \texttt{zig test}, sobre un \texttt{std.heap.ArenaAllocator} por petición HTTP, o sobre un \texttt{std.heap.FixedBufferAllocator} respaldado por un array estático en la pila sin tocar jamás el sistema operativo.
% studium:end-record eyJpZCI6ICJQQVItMDAzMyIsICJzZWN0aW9uIjogImNhcDA2LW1vbm9ncmFmaWEtemlnLWFzaWduYWRvcmVzLWNvbXB0aW1lLWludGVyb3BlcmFiaWxpZGFkLWMifQ==

\begin{tcolorbox}[title={Problema resuelto}, studiumworked, breakable]
% studium:record eyJpZCI6ICJQUkItMDAzMSIsICJzZWN0aW9uIjogImNhcDA2LW1vbm9ncmFmaWEtemlnLWFzaWduYWRvcmVzLWNvbXB0aW1lLWludGVyb3BlcmFiaWxpZGFkLWMifQ==

\noindent\textbf{Enunciado}
En un microservicio escrito en Zig, cada petición HTTP utiliza un \texttt{std.heap.FixedBufferAllocator} respaldado por un array en pila de \(B = 16384\text{ B}\) (\(16\text{ KiB}\)). Si el manejador aloja una cabecera parseada de \(512\text{ B}\) y \(N = 60\) cabeceras clave-valor de \(128\text{ B}\) cada una, calcula cuántos bytes quedan disponibles en el \texttt{FixedBufferAllocator} antes de que devuelva \texttt{error.OutOfMemory}.

\noindent\textbf{Datos}
\begin{itemize}
\item $B = 16384\,\mathrm{B}$ — Capacidad total del búfer estático en pila
\item $S_{hdr} = 512\,\mathrm{B}$ — Reserva inicial para el descriptor de petición
\item $N^{S}_{kv} = 60, 128 B$ — 60 pares clave-valor de 128 B
\end{itemize}

\noindent\textbf{Incógnita}
Magnitud objetivo solicitada en el enunciado

\noindent\textbf{Modelo físico}

\noindent\textit{Contabilidad lineal en FixedBufferAllocator de Zig}

\[
B_{\text{libre}} = 16384 - (512 + 60 \cdot 128)
\]

\noindent\textbf{Hipótesis}
\begin{itemize}
\item Evaluación determinista bajo el modelo formal del capítulo
\end{itemize}

\noindent\textbf{Desarrollo}
\begin{enumerate}
\item
Restamos los 512 B y los 60 * 128 = 7680 B a la capacidad total de 16384 B

\[
B_{\text{libre}} = 16384 - (512 + 60 \cdot 128)
\]

\[
\left(16384-\left(512+\left(60\cdot 128\right)\right)\right)
\]
\end{enumerate}

\noindent\textbf{Resultado}

\[
Resultado = 8192. Quedan exactamente 8192 bytes (8 KiB) libres en el búfer de la petición.
\]

\noindent\textbf{Interpretación}
Quedan exactamente 8192 bytes (8 KiB) libres en el búfer de la petición.

\noindent\textbf{Comprobación de los pasos}
\begin{itemize}
\item Paso 1: $8192$ — reproducido
\end{itemize}

\noindent\textbf{Comprobación de la respuesta} la aritmética del último paso reproduce la respuesta enunciada a la precisión con que está escrita. Es una reproducción aritmética; no valida el modelo físico ni sus hipótesis.
% studium:end-record eyJpZCI6ICJQUkItMDAzMSIsICJzZWN0aW9uIjogImNhcDA2LW1vbm9ncmFmaWEtemlnLWFzaWduYWRvcmVzLWNvbXB0aW1lLWludGVyb3BlcmFiaWxpZGFkLWMifQ==
\end{tcolorbox}

\begin{tcolorbox}[title={Ejercicios}, studiumworked, breakable]

\noindent\textbf{1 [PRB-0033] (Básico, Numérico)}
Una rebanada \texttt{[]const u32} en Zig de 64 bits consta de un puntero \texttt{[*]const u32} (8 B) y una longitud \texttt{usize} (8 B). ¿Cuántos bytes ocupa?

\noindent\textbf{2 [PRB-0034] (Básico, Numérico)}
En Zig, el tipo entero de anchura arbitraria \texttt{u12} se aloja en memoria ocupando el siguiente número entero de bytes potencia de dos (\texttt{@sizeOf(u12)}). ¿Cuántos bytes ocupa en memoria?

\noindent\textbf{3 [PRB-0035] (Intermedio, Desarrollo por pasos)}
Un \texttt{packed struct} de Zig contiene campos \texttt{u3}, \texttt{u5}, \texttt{u16} y \texttt{u8} empaquetados al bit exacto. ¿Cuántos bytes ocupa \texttt{(3 + 5 + 16 + 8) / 8}?

\noindent\textbf{4 [PRB-0036] (Intermedio, Diseño)}
Si una función con 4 reservas secuenciales protegidas cada una por \texttt{errdefer} falla en la cuarta reserva, ¿cuántos bloques previos son liberados automáticamente por los \texttt{errdefer} ya registrados?
\end{tcolorbox}

\chapter{Monografía de Ensamblador ARM64 (AArch64) y x86-64 (Bloques A\(-\)H): Registros, ABI, Pila y SIMD}
\label{chap:cap07-monografia-ensamblador-arm64-x86-64-abi-simd}
% studium:record eyJpZCI6ICJQQVItMDAzNyIsICJzZWN0aW9uIjogImNhcDA3LW1vbm9ncmFmaWEtZW5zYW1ibGFkb3ItYXJtNjQteDg2LTY0LWFiaS1zaW1kIn0=

\noindent\textit{Todo ingeniero de sistemas debe ser capaz de leer el código ensamblador emitido por el compilador (\texttt{clang -S -O2} u \texttt{objdump -d}), diagnosticar volcados de pila y escribir rutinas críticas vectorizadas o cambios de contexto de sistema operativo. Este capítulo ofrece el bloque completo \(A-H\) para las dos arquitecturas dominantes de 64 bits: **x86-64 (AMD64)** y **ARM64 (AArch64)**.}
% studium:end-record eyJpZCI6ICJQQVItMDAzNyIsICJzZWN0aW9uIjogImNhcDA3LW1vbm9ncmFmaWEtZW5zYW1ibGFkb3ItYXJtNjQteDg2LTY0LWFiaS1zaW1kIn0=
% studium:record eyJpZCI6ICJGSUctMDAwMiIsICJzZWN0aW9uIjogImNhcDA3LW1vbm9ncmFmaWEtZW5zYW1ibGFkb3ItYXJtNjQteDg2LTY0LWFiaS1zaW1kIn0=

\par\vspace{\baselineskip}
\noindent\begin{minipage}{\linewidth}
\centering
\begingroup
\setbox0=\hbox{\begin{tikzpicture}
\tikzset{node distance=1.4cm, >=stealth, thick, font=\small}
\node[draw=blue!70!black, fill=blue!8, rounded corners, minimum width=5.6cm, minimum height=1.8cm, align=center] (x86) {\textbf{System V AMD64 (x86-64)}\\Args 1--6: \texttt{rdi, rsi, rdx, rcx, r8, r9}\\Retorno: \texttt{rax} (\texttt{rdx})\\Callee-saved: \texttt{rbx, rbp, r12--r15}};
\node[draw=teal!70!black, fill=teal!10, rounded corners, minimum width=5.6cm, minimum height=1.8cm, align=center, right=0.8cm of x86] (arm) {\textbf{AAPCS64 (ARM64 / AArch64)}\\Args 1--8: \texttt{x0, x1, x2, x3, x4, x5, x6, x7}\\Retorno: \texttt{x0} (\texttt{x1})\\Callee-saved: \texttt{x19--x28, x29(fp), x30(lr)}};
\end{tikzpicture}}%
\count0=\ht0
\advance\count0 by \dp0
\count2=\wd0
\count4=\dimexpr\textheight/100*38\relax
\count6=\linewidth
\ifnum\count0<1
\count0=1
\fi
\ifnum\count2<1
\count2=1
\fi
\divide\count0 by 4096
\divide\count2 by 4096
\divide\count4 by 4096
\divide\count6 by 4096
\ifnum\count0<1
\count0=1
\fi
\ifnum\count2<1
\count2=1
\fi
\ifnum\count4<1
\count4=1
\fi
\ifnum\count6<1
\count6=1
\fi
\count8=\count2
\ifnum\count0>\count4
\count8=\numexpr\count2*\count4/\count0\relax
\fi
\ifnum\count8>\count6
\resizebox{\linewidth}{!}{\begin{tikzpicture}
\tikzset{node distance=1.4cm, >=stealth, thick, font=\small}
\node[draw=blue!70!black, fill=blue!8, rounded corners, minimum width=5.6cm, minimum height=1.8cm, align=center] (x86) {\textbf{System V AMD64 (x86-64)}\\Args 1--6: \texttt{rdi, rsi, rdx, rcx, r8, r9}\\Retorno: \texttt{rax} (\texttt{rdx})\\Callee-saved: \texttt{rbx, rbp, r12--r15}};
\node[draw=teal!70!black, fill=teal!10, rounded corners, minimum width=5.6cm, minimum height=1.8cm, align=center, right=0.8cm of x86] (arm) {\textbf{AAPCS64 (ARM64 / AArch64)}\\Args 1--8: \texttt{x0, x1, x2, x3, x4, x5, x6, x7}\\Retorno: \texttt{x0} (\texttt{x1})\\Callee-saved: \texttt{x19--x28, x29(fp), x30(lr)}};
\end{tikzpicture}}%
\else
\ifnum\count0>\count4
\resizebox{!}{\dimexpr\textheight/100*38\relax}{\begin{tikzpicture}
\tikzset{node distance=1.4cm, >=stealth, thick, font=\small}
\node[draw=blue!70!black, fill=blue!8, rounded corners, minimum width=5.6cm, minimum height=1.8cm, align=center] (x86) {\textbf{System V AMD64 (x86-64)}\\Args 1--6: \texttt{rdi, rsi, rdx, rcx, r8, r9}\\Retorno: \texttt{rax} (\texttt{rdx})\\Callee-saved: \texttt{rbx, rbp, r12--r15}};
\node[draw=teal!70!black, fill=teal!10, rounded corners, minimum width=5.6cm, minimum height=1.8cm, align=center, right=0.8cm of x86] (arm) {\textbf{AAPCS64 (ARM64 / AArch64)}\\Args 1--8: \texttt{x0, x1, x2, x3, x4, x5, x6, x7}\\Retorno: \texttt{x0} (\texttt{x1})\\Callee-saved: \texttt{x19--x28, x29(fp), x30(lr)}};
\end{tikzpicture}}%
\else
\begin{tikzpicture}
\tikzset{node distance=1.4cm, >=stealth, thick, font=\small}
\node[draw=blue!70!black, fill=blue!8, rounded corners, minimum width=5.6cm, minimum height=1.8cm, align=center] (x86) {\textbf{System V AMD64 (x86-64)}\\Args 1--6: \texttt{rdi, rsi, rdx, rcx, r8, r9}\\Retorno: \texttt{rax} (\texttt{rdx})\\Callee-saved: \texttt{rbx, rbp, r12--r15}};
\node[draw=teal!70!black, fill=teal!10, rounded corners, minimum width=5.6cm, minimum height=1.8cm, align=center, right=0.8cm of x86] (arm) {\textbf{AAPCS64 (ARM64 / AArch64)}\\Args 1--8: \texttt{x0, x1, x2, x3, x4, x5, x6, x7}\\Retorno: \texttt{x0} (\texttt{x1})\\Callee-saved: \texttt{x19--x28, x29(fp), x30(lr)}};
\end{tikzpicture}%
\fi
\fi
\endgroup
\par\nopagebreak
{\small \refstepcounter{studiumfigure}\label{fig:fig-abi-calling-conventions}\textbf{Figura~\thestudiumfigure.} Comparativa de paso de los primeros 6 argumentos enteros/punteros y valor de retorno entre la ABI System V x86-64 y la ABI AAPCS64 (ARM64).\par}
\end{minipage}
\par\vspace{\baselineskip}
% studium:end-record eyJpZCI6ICJGSUctMDAwMiIsICJzZWN0aW9uIjogImNhcDA3LW1vbm9ncmFmaWEtZW5zYW1ibGFkb3ItYXJtNjQteDg2LTY0LWFiaS1zaW1kIn0=

\begin{tcolorbox}[title={Consejo}, studiumtip, breakable]
% studium:record eyJpZCI6ICJQQVItMDA0MSIsICJzZWN0aW9uIjogImNhcDA3LW1vbm9ncmFmaWEtZW5zYW1ibGFkb3ItYXJtNjQteDg2LTY0LWFiaS1zaW1kIn0=

En lugar de escribir funciones enteras en ensamblador manual que dificultan las optimizaciones del compilador, prefiere funciones intrínsecas SIMD (\texttt{<arm\_neon.h>} o \texttt{<immintrin.h>}) y verifica el ensamblador generado con \texttt{clang -O3 -S}.
% studium:end-record eyJpZCI6ICJQQVItMDA0MSIsICJzZWN0aW9uIjogImNhcDA3LW1vbm9ncmFmaWEtZW5zYW1ibGFkb3ItYXJtNjQteDg2LTY0LWFiaS1zaW1kIn0=
\end{tcolorbox}

\begin{tcolorbox}[title={Definición}, studiumdef, breakable]
% studium:record eyJpZCI6ICJQQVItMDA0MCIsICJzZWN0aW9uIjogImNhcDA3LW1vbm9ncmFmaWEtZW5zYW1ibGFkb3ItYXJtNjQteDg2LTY0LWFiaS1zaW1kIn0=

Registros preservados por el llamado (Callee-saved registers): subconjunto de registros de la CPU (\texttt{rbx, rbp, r12\ensuremath{-}r15} en System V x86-64; \texttt{x19\ensuremath{-}x28, x29, x30} en AAPCS64) cuyo valor original debe ser restaurado por cualquier subrutina antes de ejecutar \texttt{ret} si dicha subrutina decide modificarlos.
% studium:end-record eyJpZCI6ICJQQVItMDA0MCIsICJzZWN0aW9uIjogImNhcDA3LW1vbm9ncmFmaWEtZW5zYW1ibGFkb3ItYXJtNjQteDg2LTY0LWFiaS1zaW1kIn0=
\end{tcolorbox}

\begin{tcolorbox}[title={Autoficha}, studiumcheck, breakable]
% studium:record eyJpZCI6ICJQQVItMDA0MiIsICJzZWN0aW9uIjogImNhcDA3LW1vbm9ncmFmaWEtZW5zYW1ibGFkb3ItYXJtNjQteDg2LTY0LWFiaS1zaW1kIn0=

¿Por qué en x86-64 al entrar en el cuerpo de una función tras ejecutar la instrucción \texttt{call} (que empuja 8 bytes de dirección de retorno en la pila), el registro \texttt{rsp} tiene un valor congruente con 8 módulo 16 y requiere restar \(8 + 16k\) bytes antes de hacer otro \texttt{call}?
% studium:end-record eyJpZCI6ICJQQVItMDA0MiIsICJzZWN0aW9uIjogImNhcDA3LW1vbm9ncmFmaWEtZW5zYW1ibGFkb3ItYXJtNjQteDg2LTY0LWFiaS1zaW1kIn0=
\end{tcolorbox}

\section{Bloques A\(-\)D: Arquitectura de registros x86-64 y ARM64 (AArch64), modos de direccionamiento y saltos}
% studium:record eyJpZCI6ICJQQVItMDAzOCIsICJzZWN0aW9uIjogImNhcDA3LW1vbm9ncmFmaWEtZW5zYW1ibGFkb3ItYXJtNjQteDg2LTY0LWFiaS1zaW1kIn0=

En **x86-64**, la CPU dispone de 16 registros de propósito general de 64 bits (\texttt{rax}, \texttt{rbx}, \texttt{rcx}, \texttt{rdx}, \texttt{rsi}, \texttt{rdi}, \texttt{rbp}, \texttt{rsp}, \texttt{r8}\(-\)\texttt{r15}) e instrucciones de longitud variable (de 1 a 15 bytes) que permiten combinar acceso a memoria y operación aritmética en una sola instrucción. En **ARM64 (AArch64)**, la arquitectura es estrictamente *load/store* con instrucciones de anchura fija de 32 bits (4 bytes) y 31 registros de propósito general de 64 bits (\texttt{x0}\(-\)\texttt{x30}, con sus mitades inferiores de 32 bits \texttt{w0}\(-\)\texttt{w30}), más el registro cero \texttt{xzr} y el puntero de pila \texttt{sp}. En ARM64, las operaciones aritméticas operan exclusivamente entre registros; la memoria solo se toca mediante instrucciones de carga y almacenamiento como \texttt{ldr}, \texttt{str}, \texttt{ldp} (cargar par de registros) y \texttt{stp} (almacenar par de registros).
% studium:end-record eyJpZCI6ICJQQVItMDAzOCIsICJzZWN0aW9uIjogImNhcDA3LW1vbm9ncmFmaWEtZW5zYW1ibGFkb3ItYXJtNjQteDg2LTY0LWFiaS1zaW1kIn0=

\section{Bloques E\(-\)H: Convenciones ABI (System V y AAPCS64), alineamiento de pila, SIMD (NEON/AVX) e inspección con objdump}
% studium:record eyJpZCI6ICJQQVItMDAzOSIsICJzZWN0aW9uIjogImNhcDA3LW1vbm9ncmFmaWEtZW5zYW1ibGFkb3ItYXJtNjQteDg2LTY0LWFiaS1zaW1kIn0=

Para que módulos compilados por distintos lenguajes (C, C++, Rust, Zig) puedan invocarse entre sí, deben respetar la **ABI** (*Application Binary Interface*) de la plataforma. Como resume la Figura \textbackslash{}ref\{fig:fig-abi-calling-conventions\}, en la ABI System V x86-64 (Linux, macOS, BSD) los primeros 6 argumentos enteros o punteros se pasan en \texttt{rdi}, \texttt{rsi}, \texttt{rdx}, \texttt{rcx}, \texttt{r8}, \texttt{r9}, mientras que en AAPCS64 (ARM64) los primeros 8 se pasan en \texttt{x0}\(-\)\texttt{x7}. En ambas arquitecturas, el puntero de pila (\texttt{rsp} o \texttt{sp}) debe mantenerse alineado a un múltiplo de **16 bytes** en el punto de llamada a cualquier función (\texttt{call} o \texttt{bl}). Compare la misma función \texttt{uint64\_t fma\_u64(uint64\_t a, uint64\_t b, uint64\_t c)} que calcula \(a \cdot b + c\) en ambas arquitecturas:
\begin{lstlisting}[style=studium,escapeinside={(*@}{@*)}]
; System V x86-64 (a=rdi, b=rsi, c=rdx -> retorno en rax):
fma_u64_x86:
    mov     rax, rdi
    imul    rax, rsi
    add     rax, rdx
    ret

; AAPCS64 ARM64 (a=x0, b=x1, c=x2 -> retorno en x0, una sola instrucci(*@{\'o}@*)n madd):
fma_u64_arm64:
    madd    x0, x0, x1, x2   ; x0 = (x0 * x1) + x2
    ret
\end{lstlisting}
En ARM64, el prólogo estándar de una función no hoja guarda el puntero de marco \texttt{x29} (\texttt{fp}) y el registro de enlace \texttt{x30} (\texttt{lr}) reservando un múltiplo de 16 bytes en una sola instrucción con pre-indexación: \texttt{stp x29, x30, [sp, \#-16]!}.
% studium:end-record eyJpZCI6ICJQQVItMDAzOSIsICJzZWN0aW9uIjogImNhcDA3LW1vbm9ncmFmaWEtZW5zYW1ibGFkb3ItYXJtNjQteDg2LTY0LWFiaS1zaW1kIn0=

\begin{tcolorbox}[title={Problema resuelto}, studiumworked, breakable]
% studium:record eyJpZCI6ICJQUkItMDAzNyIsICJzZWN0aW9uIjogImNhcDA3LW1vbm9ncmFmaWEtZW5zYW1ibGFkb3ItYXJtNjQteDg2LTY0LWFiaS1zaW1kIn0=

\noindent\textbf{Enunciado}
Una función en ensamblador ARM64 guarda en la pila el par \texttt{(x29, x30)} (\(16\text{ B}\)), dos registros preservados \texttt{(x19, x20)} (\(16\text{ B}\)) y necesita alojar un búfer local de \(36\text{ B}\). Como el hardware AArch64 exige que el puntero de pila \texttt{sp} permanezca siempre alineado a múltiplo de \(16\text{ B}\) (redondeando los \(36\text{ B}\) locales hacia arriba a \(48\text{ B}\)), calcula cuántos bytes totales resta al registro \texttt{sp} el prólogo de la función.

\noindent\textbf{Datos}
\begin{itemize}
\item $S_{fplr} = 16\,\mathrm{B}$ — Par de registros (x29, x30)
\item $S_{callee} = 16\,\mathrm{B}$ — Par de registros preservados (x19, x20)
\item $S_{\mathcal{lo} aligned} = 48\,\mathrm{B}$ — 36 B redondeados al siguiente múltiplo de 16 B
\end{itemize}

\noindent\textbf{Incógnita}
Magnitud objetivo solicitada en el enunciado

\noindent\textbf{Modelo físico}

\noindent\textit{Cálculo del tamaño de marco de pila alineado a 16 bytes en AAPCS64}

\[
S_{\text{frame}} = 16 + 16 + 48
\]

\noindent\textbf{Hipótesis}
\begin{itemize}
\item Evaluación determinista bajo el modelo formal del capítulo
\end{itemize}

\noindent\textbf{Desarrollo}
\begin{enumerate}
\item
Sumamos los tres bloques alineados a 16 bytes

\[
S_{\text{frame}} = 16 + 16 + 48
\]

\[
\left(\left(16+16\right)+48\right)
\]
\end{enumerate}

\noindent\textbf{Resultado}

\[
Resultado = 80. El prólogo reserva exactamente 80 bytes (múltiplo exacto de 16: 5 * 16 = 80 B) en la pila.
\]

\noindent\textbf{Interpretación}
El prólogo reserva exactamente 80 bytes (múltiplo exacto de 16: 5 * 16 = 80 B) en la pila.

\noindent\textbf{Comprobación de los pasos}
\begin{itemize}
\item Paso 1: $80$ — reproducido
\end{itemize}

\noindent\textbf{Comprobación de la respuesta} la aritmética del último paso reproduce la respuesta enunciada a la precisión con que está escrita. Es una reproducción aritmética; no valida el modelo físico ni sus hipótesis.
% studium:end-record eyJpZCI6ICJQUkItMDAzNyIsICJzZWN0aW9uIjogImNhcDA3LW1vbm9ncmFmaWEtZW5zYW1ibGFkb3ItYXJtNjQteDg2LTY0LWFiaS1zaW1kIn0=
\end{tcolorbox}

\begin{tcolorbox}[title={Ejercicios}, studiumworked, breakable]

\noindent\textbf{1 [PRB-0039] (Básico, Numérico)}
En ARM64, cada instrucción máquina tiene una anchura fija de 32 bits (4 bytes). ¿Cuántos bytes ocupan en la sección \texttt{.text} una secuencia de 250 instrucciones?

\noindent\textbf{2 [PRB-0040] (Básico, Numérico)}
Un registro vectorial NEON de 128 bits en ARM64 (\texttt{v0.4s}) procesa enteros o flotantes de 32 bits en paralelo. ¿Cuántos elementos de 32 bits opera por instrucción?

\noindent\textbf{3 [PRB-0041] (Intermedio, Desarrollo por pasos)}
Una función en System V x86-64 recibe 9 argumentos enteros de 64 bits. Sabiendo que los primeros 6 van en registros y el resto se pasa por la pila (8 B cada uno), ¿cuántos bytes ocupan en la pila los argumentos restantes?

\noindent\textbf{4 [PRB-0042] (Intermedio, Diseño)}
Un bucle procesa un array de 4096 flotantes \texttt{f32} utilizando instrucciones AVX-512 (16 flotantes \texttt{f32} por registro de 512 bits). ¿Cuántas iteraciones vectoriales ejecuta el bucle?
\end{tcolorbox}

\chapter{Monografías de Ada/SPARK (Sistemas Críticos) y Fortran Moderno (Cómputo Científico HPC) (Bloques A\(-\)H)}
\label{chap:cap08-monografias-ada-spark-tiempo-real-y-fortran-cientifico}
% studium:record eyJpZCI6ICJQQVItMDA0MyIsICJzZWN0aW9uIjogImNhcDA4LW1vbm9ncmFmaWFzLWFkYS1zcGFyay10aWVtcG8tcmVhbC15LWZvcnRyYW4tY2llbnRpZmljbyJ9

\noindent\textit{Este capítulo cierra el Volumen IV con las monografías completas de dos lenguajes insustituibles en sus dominios de alta exigencia: **Ada 2022 / SPARK** en aviónica (DO-178C Nivel A), ferrocarriles y espacio, y **Fortran 2018/2023** en supercomputación (HPC), dinámica de fluidos computacional (CFD) y predicción meteorológica.}
% studium:end-record eyJpZCI6ICJQQVItMDA0MyIsICJzZWN0aW9uIjogImNhcDA4LW1vbm9ncmFmaWFzLWFkYS1zcGFyay10aWVtcG8tcmVhbC15LWZvcnRyYW4tY2llbnRpZmljbyJ9

\begin{tcolorbox}[title={Consejo}, studiumtip, breakable]
% studium:record eyJpZCI6ICJQQVItMDA0NyIsICJzZWN0aW9uIjogImNhcDA4LW1vbm9ncmFmaWFzLWFkYS1zcGFyay10aWVtcG8tcmVhbC15LWZvcnRyYW4tY2llbnRpZmljbyJ9

En Fortran moderno, comienza siempre todo módulo y programa con \texttt{implicit none} e importa \texttt{real64} (\texttt{dp}) e \texttt{int64} desde el módulo estándar \texttt{iso\_fortran\_env}; al anidar bucles sobre \texttt{A(i, j)}, sitúa siempre el índice \texttt{i} en el bucle más interno y \texttt{j} en el externo para recorrer la memoria secuencialmente.
% studium:end-record eyJpZCI6ICJQQVItMDA0NyIsICJzZWN0aW9uIjogImNhcDA4LW1vbm9ncmFmaWFzLWFkYS1zcGFyay10aWVtcG8tcmVhbC15LWZvcnRyYW4tY2llbnRpZmljbyJ9
\end{tcolorbox}

\begin{tcolorbox}[title={Definición}, studiumdef, breakable]
% studium:record eyJpZCI6ICJQQVItMDA0NiIsICJzZWN0aW9uIjogImNhcDA4LW1vbm9ncmFmaWFzLWFkYS1zcGFyay10aWVtcG8tcmVhbC15LWZvcnRyYW4tY2llbnRpZmljbyJ9

Orden mayor por columnas (Column-Major Order en Fortran): disposición lineal en memoria de un array bidimensional \(A(1..M, 1..N)\) según la cual el elemento \(A(i, j)\) (con índices basados en 1) se ubica en el desplazamiento lineal \((j - 1) \cdot M + (i - 1)\), haciendo contiguos los elementos de una misma columna.
% studium:end-record eyJpZCI6ICJQQVItMDA0NiIsICJzZWN0aW9uIjogImNhcDA4LW1vbm9ncmFmaWFzLWFkYS1zcGFyay10aWVtcG8tcmVhbC15LWZvcnRyYW4tY2llbnRpZmljbyJ9
\end{tcolorbox}

\begin{tcolorbox}[title={Autoficha}, studiumcheck, breakable]
% studium:record eyJpZCI6ICJQQVItMDA0OCIsICJzZWN0aW9uIjogImNhcDA4LW1vbm9ncmFmaWFzLWFkYS1zcGFyay10aWVtcG8tcmVhbC15LWZvcnRyYW4tY2llbnRpZmljbyJ9

¿Por qué intercambiar el orden de los bucles \texttt{do i = 1, M; do j = 1, N} sobre \texttt{A(i, j)} en Fortran degrada severamente la tasa de aciertos de caché L1 respecto a \texttt{do j = 1, N; do i = 1, M}?
% studium:end-record eyJpZCI6ICJQQVItMDA0OCIsICJzZWN0aW9uIjogImNhcDA4LW1vbm9ncmFmaWFzLWFkYS1zcGFyay10aWVtcG8tcmVhbC15LWZvcnRyYW4tY2llbnRpZmljbyJ9
\end{tcolorbox}

\section{Monografía de Ada 2022 y SPARK (Bloques A\(-\)H): Tipado fuerte de dominio, contratos y sistemas críticos}
% studium:record eyJpZCI6ICJQQVItMDA0NCIsICJzZWN0aW9uIjogImNhcDA4LW1vbm9ncmFmaWFzLWFkYS1zcGFyay10aWVtcG8tcmVhbC15LWZvcnRyYW4tY2llbnRpZmljbyJ9

Encargado por el Departamento de Defensa de EE. UU. (Jean Ichbiah, 1980; revisiones ISO Ada 95, 2005, 2012 y 2022), **Ada** fue diseñado para sistemas embebidos de misión crítica donde un defecto de software puede costar vidas humanas. En Ada, el ingeniero no declara variables como meros \texttt{int}, sino que define tipos nominalmente distintos con su rango físico exacto: \texttt{type Altitude\_Meters is range -500 .. 50\_000;} y \texttt{type Velocity\_Knots is range 0 .. 2\_500;}. Sumar una altitud con una velocidad es un error de compilación inmediato, y asignar un valor fuera de rango lanza \texttt{Constraint\_Error}. En el subconjunto formalmente verificable **SPARK**, los aspectos \texttt{Pre}, \texttt{Post} y \texttt{Global} se traducen en obligaciones de prueba matemática descargadas por solucionadores SMT (\texttt{Z3}, \texttt{Alt-Ergo}) mediante \texttt{gnatprove}, demostrando la ausencia total de desbordamientos o divisiones por cero antes de compilar:
\begin{lstlisting}[style=studium]
package Flight_Envelope with SPARK_Mode is
   type Altitude_M is range 0 .. 20_000;
   function Safe_Climb (Current : Altitude_M; Delta_H : Altitude_M) return Altitude_M
     with Pre  => Integer (Current) + Integer (Delta_H) <= 20_000,
          Post => Safe_Climb'Result = Current + Delta_H;
end Flight_Envelope;
\end{lstlisting}
Para concurrencia de tiempo real duro, el perfil **Ravenscar** restringe las tareas (\texttt{task}) y los objetos protegidos (\texttt{protected object}) al protocolo de techo de prioridad inmediata (*Immediate Ceiling Priority Protocol*), garantizando ausencia matemática de interbloqueos (*deadlocks*) y acotación estática del tiempo de respuesta.
% studium:end-record eyJpZCI6ICJQQVItMDA0NCIsICJzZWN0aW9uIjogImNhcDA4LW1vbm9ncmFmaWFzLWFkYS1zcGFyay10aWVtcG8tcmVhbC15LWZvcnRyYW4tY2llbnRpZmljbyJ9

\section{Monografía de Fortran 2018/2023 (Bloques A\(-\)H): Arrays multidimensionales, vectorización y Coarrays}
% studium:record eyJpZCI6ICJQQVItMDA0NSIsICJzZWN0aW9uIjogImNhcDA4LW1vbm9ncmFmaWFzLWFkYS1zcGFyay10aWVtcG8tcmVhbC15LWZvcnRyYW4tY2llbnRpZmljbyJ9

Por su parte, **Fortran** (John Backus, IBM, 1957; modernizado radicalmente en Fortran 90/95, 2003, 2008, 2018 y 2023) sigue siendo el estándar de oro en supercomputación científica por dos razones técnicas: primero, los arrays multidimensionales son ciudadanos de primera clase almacenados en **orden mayor por columnas** (*column-major*, donde \texttt{A(i, j)} es contiguo en memoria al variar el primer índice \texttt{i}); segundo, el estándar Fortran garantiza por defecto que los argumentos \texttt{intent(in)} e \texttt{intent(inout)} de una subrutina \texttt{pure} **no presentan alias entre sí** (a menos que tengan el atributo \texttt{pointer} o \texttt{target}), lo que permite a compiladores como \texttt{gfortran}, \texttt{ifx} o \texttt{nvfortran} vectorizar bucles matriciales con una agresividad superior a la de C sin \texttt{restrict}:
\begin{lstlisting}[style=studium]
module numerics_mod
  use, intrinsic :: iso_fortran_env, only: dp => real64
  implicit none
contains
  pure subroutine laplacian_1d(u, d2u, inv_dx2)
    real(dp), intent(in)  :: u(:), inv_dx2
    real(dp), intent(out) :: d2u(:)
    integer :: n
    n = size(u)
    d2u(1) = 0.0_dp
    d2u(2:n-1) = (u(1:n-2) - 2.0_dp * u(2:n-1) + u(3:n)) * inv_dx2
    d2u(n) = 0.0_dp
  end subroutine laplacian_1d
end module numerics_mod
\end{lstlisting}
Además, Fortran 2008/2018 incorpora paralelismo de memoria distribuida directamente en la sintaxis del lenguaje mediante **Coarrays** (\texttt{real(dp) :: field(1000)[*]}), donde la notación entre corchetes \texttt{[imagen]} realiza transferencias RDMA de un solo lado entre nodos de un clúster sin llamadas verbosas a MPI.
% studium:end-record eyJpZCI6ICJQQVItMDA0NSIsICJzZWN0aW9uIjogImNhcDA4LW1vbm9ncmFmaWFzLWFkYS1zcGFyay10aWVtcG8tcmVhbC15LWZvcnRyYW4tY2llbnRpZmljbyJ9

\begin{tcolorbox}[title={Problema resuelto}, studiumworked, breakable]
% studium:record eyJpZCI6ICJQUkItMDA0MyIsICJzZWN0aW9uIjogImNhcDA4LW1vbm9ncmFmaWFzLWFkYS1zcGFyay10aWVtcG8tcmVhbC15LWZvcnRyYW4tY2llbnRpZmljbyJ9

\noindent\textbf{Enunciado}
En un programa Fortran 2018, una matriz \texttt{real(real64) :: A(100, 200)} (\(M = 100\) filas, \(N = 200\) columnas, \(S = 8\text{ B}\) por elemento) está almacenada en orden mayor por columnas a partir del desplazamiento base \(0\). Calcula el desplazamiento exacto en bytes del elemento \texttt{A(25, 4)} usando la fórmula \(8 \cdot ((4 - 1) \cdot 100 + (25 - 1))\).

\noindent\textbf{Datos}
\begin{itemize}
\item $M = 100$ — Número de filas por columna
\item $i^{j} = 25, 4$ — Fila 25, columna 4 (base 1)
\item $S = 8\,\mathrm{B}$ — Tamaño de real(real64)
\end{itemize}

\noindent\textbf{Incógnita}
Magnitud objetivo solicitada en el enunciado

\noindent\textbf{Modelo físico}

\noindent\textit{Direccionamiento Column-Major en Fortran}

\[
\text{Offset} = 8 \cdot ((4 - 1) \cdot 100 + (25 - 1))
\]

\noindent\textbf{Hipótesis}
\begin{itemize}
\item Evaluación determinista bajo el modelo formal del capítulo
\end{itemize}

\noindent\textbf{Desarrollo}
\begin{enumerate}
\item
Calculamos el índice lineal (3 * 100 + 24 = 324) y multiplicamos por 8 bytes

\[
\text{Offset} = 8 \cdot ((4 - 1) \cdot 100 + (25 - 1))
\]

\[
\left(8\cdot \left(\left(\left(4-1\right)\cdot 100\right)+\left(25-1\right)\right)\right)
\]
\end{enumerate}

\noindent\textbf{Resultado}

\[
Resultado = 2592. El elemento A(25, 4) se encuentra exactamente a 2592 bytes del inicio de la matriz.
\]

\noindent\textbf{Interpretación}
El elemento A(25, 4) se encuentra exactamente a 2592 bytes del inicio de la matriz.

\noindent\textbf{Comprobación de los pasos}
\begin{itemize}
\item Paso 1: $2592$ — reproducido
\end{itemize}

\noindent\textbf{Comprobación de la respuesta} la aritmética del último paso reproduce la respuesta enunciada a la precisión con que está escrita. Es una reproducción aritmética; no valida el modelo físico ni sus hipótesis.
% studium:end-record eyJpZCI6ICJQUkItMDA0MyIsICJzZWN0aW9uIjogImNhcDA4LW1vbm9ncmFmaWFzLWFkYS1zcGFyay10aWVtcG8tcmVhbC15LWZvcnRyYW4tY2llbnRpZmljbyJ9
\end{tcolorbox}

\begin{tcolorbox}[title={Ejercicios}, studiumworked, breakable]

\noindent\textbf{1 [PRB-0045] (Básico, Numérico)}
En Ada, el tipo \texttt{type Percent is range 0 .. 100;} comprende desde 0 hasta 100 ambos inclusive. ¿Cuántos valores distintos pertenecen al tipo (\texttt{Percent'Last - Percent'First + 1})?

\noindent\textbf{2 [PRB-0046] (Básico, Numérico)}
En Fortran, una rebanada vectorial \texttt{u(2:n-1)} para \(n = 1002\) evalúa los puntos interiores de una malla 1D. ¿Cuántos puntos interiores actualiza la sentencia?

\noindent\textbf{3 [PRB-0047] (Intermedio, Desarrollo por pasos)}
Un programa Coarray Fortran se ejecuta en 16 imágenes (\texttt{num\_images() = 16}) y cada imagen aloja una malla local de \(50000\) flotantes \texttt{real64} (8 B). ¿Cuántos bytes suman las mallas de las 16 imágenes?

\noindent\textbf{4 [PRB-0048] (Intermedio, Diseño)}
En un sistema de tiempo real Ada Ravenscar con periodo de marco mayor \(T = 100\text{ ms}\), tres tareas periódicas consumen \(15\text{ ms}\), \(22\text{ ms}\) y \(18\text{ ms}\) por cada \(100\text{ ms}\). ¿Cuál es el porcentaje entero de utilización total de CPU?
\end{tcolorbox}
\appendix
\chapter{Notación}

\begingroup
\setlength{\LTpre}{0pt}
\setlength{\LTpost}{0pt}
\begin{longtable}{@{}>{\centering\arraybackslash}p{0.17\linewidth}>{\raggedright\arraybackslash}p{0.55\linewidth}>{\raggedright\arraybackslash}p{0.2\linewidth}@{}}
\textbf{Símbolo} & \textbf{Significado} & \textbf{Unidades} \\
\hline
\endfirsthead
\textbf{Símbolo} & \textbf{Significado} & \textbf{Unidades} \\
\hline
\endhead
$\text{alignof}(T)$ & Requisito de alineamiento de dirección en C11/C23 y C++11/C++20 & $\mathrm{B}$ \\
$\text{sizeof}(T)$ & Tamaño en bytes de un objeto o tipo completo en C/C++ & $\mathrm{B}$ \\
$k + p$ & Aritmética de punteros escalada por k * sizeof(*p) bytes & $\mathrm{dirección}$ \\
\end{longtable}
\endgroup
\chapter{Hoja de fórmulas}

\noindent Esta parte del borrador aún no tiene material.
\chapter{Soluciones}

\section{Monografía de C (I \(-\) Bloques A, B y C): Historia, modelo de traducción ISO C17/C23, toolchain y nivel principiante}

\noindent\textbf{Problema resuelto [PRB-0001]}

En una arquitectura de 32 bits para enteros \texttt{uint32\_t} (\(w = 32\), módulo \(2^{32} = 4294967296\)), un entero con signo con valor \(-100\) se convierte implícitamente a \texttt{uint32\_t} durante una comparación. Calcula el valor numérico exacto que toma tras la conversión según la regla \(2^{32} - 100\).

\noindent\textbf{Respuesta}

\[
Resultado = 4294967196. El valor resultante en uint32_t es 4 294 967 196, razón por la cual `-Wsign-conversion` es indispensable.
\]

\noindent\textbf{Comprobación de la respuesta} la aritmética del último paso reproduce la respuesta enunciada a la precisión con que está escrita. Es una reproducción aritmética; no valida el modelo físico ni sus hipótesis.

\noindent\textbf{Ejercicio 1 [PRB-0003]}

Calcula el valor en \texttt{uint16\_t} (\(w = 16\)) resultante de convertir \(-1\) a sin signo (\(2^{16} - 1\)).

\noindent\textbf{Respuesta} 65535

\noindent\textbf{Ejercicio 2 [PRB-0004]}

Un registro de control de 8 bits activa los bits 0, 3 y 7 (\$\(2^0 + 2^3 + 2^7\)\$). ¿Cuál es el valor decimal de la máscara \texttt{uint8\_t}?

\noindent\textbf{Respuesta} 137

\noindent\textbf{Ejercicio 3 [PRB-0005]}

¿Cuántos bytes ocupa en total una estructura con 4 campos \texttt{uint32\_t} y 2 campos \texttt{uint64\_t} perfectamente alineados (\(4 \times 4 + 2 \times 8\))?

\noindent\textbf{Respuesta} 32

\noindent\textbf{Ejercicio 4 [PRB-0006]}

Si compilamos un proyecto en C formado por 18 unidades de traducción \texttt{.c} y solo modificamos un fichero \texttt{.c}, ¿cuántas unidades de traducción evita recompilar \texttt{make}/\texttt{ninja}?

\noindent\textbf{Respuesta} 17

\subsection*{Cálculos reproducidos}

\(\left(\left(2\right)^{32}-100\right) = 4294967196\)

\section{Monografía de C (II \(-\) Bloques D y E): Punteros, aritmética de direcciones, memoria dinámica, structs y strict aliasing}

\noindent\textbf{Problema resuelto [PRB-0007]}

En el asignador por región \texttt{projects/c-arena-allocator}, el desplazamiento actual del búfer es \texttt{offset = 19} bytes y se solicita reservar una estructura de \(40\text{ B}\) con requisito de alineamiento \texttt{align = 8} bytes. El asignador avanza primero \texttt{offset} al siguiente múltiplo de 8 (sumando \(5\text{ B}\) de relleno hasta \(24\)) y luego suma los \(40\text{ B}\) del objeto. Calcula el nuevo valor exacto de \texttt{offset} tras la asignación.

\noindent\textbf{Respuesta}

\[
Resultado = 64. El desplazamiento final de la arena queda en 64 bytes.
\]

\noindent\textbf{Comprobación de la respuesta} la aritmética del último paso reproduce la respuesta enunciada a la precisión con que está escrita. Es una reproducción aritmética; no valida el modelo físico ni sus hipótesis.

\noindent\textbf{Ejercicio 1 [PRB-0009]}

Si un puntero \texttt{uint64\_t *p} apunta a la dirección decimal 4096, ¿a qué dirección decimal apunta \texttt{p + 5} sabiendo que \texttt{sizeof(uint64\_t) = 8}?

\noindent\textbf{Respuesta} 4136

\noindent\textbf{Ejercicio 2 [PRB-0010]}

Una cadena en C contiene 31 caracteres visibles más su byte terminador \texttt{'\textbackslash{}0'}. ¿Cuántos bytes mínimos debe reservar \texttt{malloc}?

\noindent\textbf{Respuesta} 32

\noindent\textbf{Ejercicio 3 [PRB-0011]}

En una arena con capacidad total de 4096 bytes, se reservan 12 nodos de 64 bytes ya alineados. ¿Cuántos bytes libres quedan en la arena?

\noindent\textbf{Respuesta} 3328

\noindent\textbf{Ejercicio 4 [PRB-0012]}

Si un paquete binario de telemetría tiene 8 bytes de cabecera y transporta 16 muestras \texttt{int16\_t} (2 bytes cada una), ¿cuántos bytes mide el paquete completo?

\noindent\textbf{Respuesta} 40

\subsection*{Cálculos reproducidos}

\(\left(\left(19+5\right)+40\right) = 64\)

\section{Monografía de C (III \(-\) Bloques F, G y H): Concurrencia C11/POSIX, Sanitizers, Ecosistema y Proyectos}

\noindent\textbf{Problema resuelto [PRB-0013]}

Cuando un programa en C se compila con \texttt{-fsanitize=address} (AddressSanitizer), cada \(8\text{ B}\) de memoria de aplicación requieren \(1\text{ B}\) de memoria sombra (*shadow memory*), más zonas rojas (*redzones*) de \(32\text{ B}\) alrededor de cada asignación. Si una aplicación reserva \(N = 1000\) bloques de \(64\text{ B}\) cada uno (con \(32\text{ B}\) de zona roja por bloque), calcula cuántos bytes suman los bloques más sus zonas rojas más la memoria sombra correspondiente (\$1/8\$ del total).

\noindent\textbf{Respuesta}

\[
Resultado = 108000. El conjunto instrumentado ocupa 108 000 bytes (96 000 B de bloques con redzones y 12 000 B de memoria sombra).
\]

\noindent\textbf{Comprobación de la respuesta} la aritmética del último paso reproduce la respuesta enunciada a la precisión con que está escrita. Es una reproducción aritmética; no valida el modelo físico ni sus hipótesis.

\noindent\textbf{Ejercicio 1 [PRB-0015]}

En AddressSanitizer, ¿cuántos bytes de memoria sombra corresponden a 64 000 bytes de memoria de aplicación (\$64000 / 8\$)?

\noindent\textbf{Respuesta} 8000

\noindent\textbf{Ejercicio 2 [PRB-0016]}

En \texttt{projects/c-arena-allocator}, la cabecera de telemetría suma los campos \texttt{magic} (10), \texttt{version} (2), \texttt{flags} (3) y \texttt{payload\_len} (25). ¿Cuál es el valor del checksum esperado?

\noindent\textbf{Respuesta} 40

\noindent\textbf{Ejercicio 3 [PRB-0017]}

Una función en C adquiere secuencialmente 5 recursos y falla al intentar adquirir el cuarto. ¿Cuántos recursos ya adquiridos debe liberar la cascada \texttt{goto cleanup} antes de retornar?

\noindent\textbf{Respuesta} 3

\noindent\textbf{Ejercicio 4 [PRB-0018]}

Un anillo de descriptores E/S tiene capacidad para 256 entradas de 16 bytes cada una. ¿Cuántos bytes mide el anillo completo mapeado en memoria?

\noindent\textbf{Respuesta} 4096

\subsection*{Cálculos reproducidos}

\(\frac{\left(\left(1000\cdot \left(64+32\right)\right)\cdot 9\right)}{8} = 108000\)

\section{Monografía de C++ (I \(-\) Bloques A\(-\)D): De C++98 a C++20/23, RAII, Semántica de Movimiento y Punteros Inteligentes}

\noindent\textbf{Problema resuelto [PRB-0019]}

En una implementación estándar de C++ de 64 bits donde \texttt{sizeof(std::string) = 24} bytes y la optimización de cadena corta (*Small String Optimization*, SSO) almacena en línea cadenas de hasta \(22\text{ B}\) sin alojar en el montículo, un programa crea un \texttt{std::vector<std::string>} con \(N = 5000\) identificadores cortos de \(12\text{ B}\) cada uno. Calcula cuántos bytes totales ocupan los 5000 objetos \texttt{std::string} dentro del búfer contiguo del vector (sabiendo que realizan \(0\) alojamientos adicionales fuera del vector).

\noindent\textbf{Respuesta}

\[
Resultado = 120000. Los 5000 identificadores se almacenan contiguamente en 120 000 bytes con cero llamadas individuales al montículo por cadena.
\]

\noindent\textbf{Comprobación de la respuesta} la aritmética del último paso reproduce la respuesta enunciada a la precisión con que está escrita. Es una reproducción aritmética; no valida el modelo físico ni sus hipótesis.

\noindent\textbf{Ejercicio 1 [PRB-0021]}

¿Cuántos bytes ocupa \texttt{std::unique\_ptr<double>} con deleción por defecto en una arquitectura de 64 bits?

\noindent\textbf{Respuesta} 8

\noindent\textbf{Ejercicio 2 [PRB-0022]}

Un \texttt{std::span<const int>} en C++20 con extensión dinámica almacena un puntero (8 B) y una longitud \texttt{size\_t} (8 B). ¿Cuántos bytes ocupa en la pila?

\noindent\textbf{Respuesta} 16

\noindent\textbf{Ejercicio 3 [PRB-0023]}

Un \texttt{std::vector<int>} duplica su capacidad geométricamente desde 16 hasta 1024 elementos (\texttt{sizeof(int) = 4}). ¿Cuántos bytes reserva el búfer final de 1024 enteros?

\noindent\textbf{Respuesta} 4096

\noindent\textbf{Ejercicio 4 [PRB-0024]}

Un bloque de control de \texttt{std::make\_shared<T>} combina en una única reserva de montículo los dos contadores atómicos (16 B) más el objeto \texttt{T} de 48 B. ¿Cuántos bytes mide esa única reserva?

\noindent\textbf{Respuesta} 64

\subsection*{Cálculos reproducidos}

\(\left(5000\cdot 24\right) = 120000\)

\section{Monografía de C++ (II \(-\) Bloques E\(-\)H): Plantillas, Concepts C++20, std::ranges, Concurrencia, Corrutinas y Proyectos}

\noindent\textbf{Problema resuelto [PRB-0025]}

En el proyecto \texttt{projects/cpp-task-scheduler}, se encolan cuatro tareas \texttt{(id, prioridad, base\_cost)}: \texttt{(1, 2, 10)}, \texttt{(2, 9, 10)}, \texttt{(3, 5, 20)} y \texttt{(4, 9, 15)}. Si se ejecuta \texttt{execute\_filtered(5)} (que filtra las tareas con \texttt{priority >= 5} y suma \texttt{base\_cost * priority} de cada una), calcula el resultado exacto devuelto.

\noindent\textbf{Respuesta}

\[
Resultado = 325. El planificador devuelve exactamente 325, coincidiendo con el test ejecutable de `projects/cpp-task-scheduler`.
\]

\noindent\textbf{Comprobación de la respuesta} la aritmética del último paso reproduce la respuesta enunciada a la precisión con que está escrita. Es una reproducción aritmética; no valida el modelo físico ni sus hipótesis.

\noindent\textbf{Ejercicio 1 [PRB-0027]}

En \texttt{cpp-task-scheduler}, ¿qué valor devuelve \texttt{ComputeTask\{10, 7, 12\}.execute()} (\(12 \times 7\))?

\noindent\textbf{Respuesta} 84

\noindent\textbf{Ejercicio 2 [PRB-0028]}

Una plantilla variádica \texttt{sum\_all(10, 20, 30, 40)} utiliza una expresión pliegue (*fold expression* \texttt{(args + ...)}). ¿Cuál es el resultado evaluado en \texttt{constexpr}?

\noindent\textbf{Respuesta} 100

\noindent\textbf{Ejercicio 3 [PRB-0029]}

Un \texttt{std::variant<int, double, std::array<char, 32>>} almacena la alternativa mayor (32 B) más un índice alineado a 8 B. ¿Cuántos bytes ocupa la instancia?

\noindent\textbf{Respuesta} 40

\noindent\textbf{Ejercicio 4 [PRB-0030]}

Si una vista \texttt{std::views::iota(1, 1001) | std::views::filter(...) | std::views::take(5)} detiene la evaluación al encontrar los primeros 5 pares (\$2, 4, 6, 8, 10\$), ¿cuántos enteros del rango \$[1, 1000]\$ llegaron a inspeccionarse?

\noindent\textbf{Respuesta} 10

\subsection*{Cálculos reproducidos}

\(\left(\left(\left(9\cdot 10\right)+\left(9\cdot 15\right)\right)+\left(5\cdot 20\right)\right) = 325\)

\section{Monografía de Zig (Bloques A\(-\)H): Asignadores explícitos, metaprogramación comptime, uniones de error y toolchain C}

\noindent\textbf{Problema resuelto [PRB-0031]}

En un microservicio escrito en Zig, cada petición HTTP utiliza un \texttt{std.heap.FixedBufferAllocator} respaldado por un array en pila de \(B = 16384\text{ B}\) (\(16\text{ KiB}\)). Si el manejador aloja una cabecera parseada de \(512\text{ B}\) y \(N = 60\) cabeceras clave-valor de \(128\text{ B}\) cada una, calcula cuántos bytes quedan disponibles en el \texttt{FixedBufferAllocator} antes de que devuelva \texttt{error.OutOfMemory}.

\noindent\textbf{Respuesta}

\[
Resultado = 8192. Quedan exactamente 8192 bytes (8 KiB) libres en el búfer de la petición.
\]

\noindent\textbf{Comprobación de la respuesta} la aritmética del último paso reproduce la respuesta enunciada a la precisión con que está escrita. Es una reproducción aritmética; no valida el modelo físico ni sus hipótesis.

\noindent\textbf{Ejercicio 1 [PRB-0033]}

Una rebanada \texttt{[]const u32} en Zig de 64 bits consta de un puntero \texttt{[*]const u32} (8 B) y una longitud \texttt{usize} (8 B). ¿Cuántos bytes ocupa?

\noindent\textbf{Respuesta} 16

\noindent\textbf{Ejercicio 2 [PRB-0034]}

En Zig, el tipo entero de anchura arbitraria \texttt{u12} se aloja en memoria ocupando el siguiente número entero de bytes potencia de dos (\texttt{@sizeOf(u12)}). ¿Cuántos bytes ocupa en memoria?

\noindent\textbf{Respuesta} 2

\noindent\textbf{Ejercicio 3 [PRB-0035]}

Un \texttt{packed struct} de Zig contiene campos \texttt{u3}, \texttt{u5}, \texttt{u16} y \texttt{u8} empaquetados al bit exacto. ¿Cuántos bytes ocupa \texttt{(3 + 5 + 16 + 8) / 8}?

\noindent\textbf{Respuesta} 4

\noindent\textbf{Ejercicio 4 [PRB-0036]}

Si una función con 4 reservas secuenciales protegidas cada una por \texttt{errdefer} falla en la cuarta reserva, ¿cuántos bloques previos son liberados automáticamente por los \texttt{errdefer} ya registrados?

\noindent\textbf{Respuesta} 3

\subsection*{Cálculos reproducidos}

\(\left(16384-\left(512+\left(60\cdot 128\right)\right)\right) = 8192\)

\section{Monografía de Ensamblador ARM64 (AArch64) y x86-64 (Bloques A\(-\)H): Registros, ABI, Pila y SIMD}

\noindent\textbf{Problema resuelto [PRB-0037]}

Una función en ensamblador ARM64 guarda en la pila el par \texttt{(x29, x30)} (\(16\text{ B}\)), dos registros preservados \texttt{(x19, x20)} (\(16\text{ B}\)) y necesita alojar un búfer local de \(36\text{ B}\). Como el hardware AArch64 exige que el puntero de pila \texttt{sp} permanezca siempre alineado a múltiplo de \(16\text{ B}\) (redondeando los \(36\text{ B}\) locales hacia arriba a \(48\text{ B}\)), calcula cuántos bytes totales resta al registro \texttt{sp} el prólogo de la función.

\noindent\textbf{Respuesta}

\[
Resultado = 80. El prólogo reserva exactamente 80 bytes (múltiplo exacto de 16: 5 * 16 = 80 B) en la pila.
\]

\noindent\textbf{Comprobación de la respuesta} la aritmética del último paso reproduce la respuesta enunciada a la precisión con que está escrita. Es una reproducción aritmética; no valida el modelo físico ni sus hipótesis.

\noindent\textbf{Ejercicio 1 [PRB-0039]}

En ARM64, cada instrucción máquina tiene una anchura fija de 32 bits (4 bytes). ¿Cuántos bytes ocupan en la sección \texttt{.text} una secuencia de 250 instrucciones?

\noindent\textbf{Respuesta} 1000

\noindent\textbf{Ejercicio 2 [PRB-0040]}

Un registro vectorial NEON de 128 bits en ARM64 (\texttt{v0.4s}) procesa enteros o flotantes de 32 bits en paralelo. ¿Cuántos elementos de 32 bits opera por instrucción?

\noindent\textbf{Respuesta} 4

\noindent\textbf{Ejercicio 3 [PRB-0041]}

Una función en System V x86-64 recibe 9 argumentos enteros de 64 bits. Sabiendo que los primeros 6 van en registros y el resto se pasa por la pila (8 B cada uno), ¿cuántos bytes ocupan en la pila los argumentos restantes?

\noindent\textbf{Respuesta} 24

\noindent\textbf{Ejercicio 4 [PRB-0042]}

Un bucle procesa un array de 4096 flotantes \texttt{f32} utilizando instrucciones AVX-512 (16 flotantes \texttt{f32} por registro de 512 bits). ¿Cuántas iteraciones vectoriales ejecuta el bucle?

\noindent\textbf{Respuesta} 256

\subsection*{Cálculos reproducidos}

\(\left(\left(16+16\right)+48\right) = 80\)

\section{Monografías de Ada/SPARK (Sistemas Críticos) y Fortran Moderno (Cómputo Científico HPC) (Bloques A\(-\)H)}

\noindent\textbf{Problema resuelto [PRB-0043]}

En un programa Fortran 2018, una matriz \texttt{real(real64) :: A(100, 200)} (\(M = 100\) filas, \(N = 200\) columnas, \(S = 8\text{ B}\) por elemento) está almacenada en orden mayor por columnas a partir del desplazamiento base \(0\). Calcula el desplazamiento exacto en bytes del elemento \texttt{A(25, 4)} usando la fórmula \(8 \cdot ((4 - 1) \cdot 100 + (25 - 1))\).

\noindent\textbf{Respuesta}

\[
Resultado = 2592. El elemento A(25, 4) se encuentra exactamente a 2592 bytes del inicio de la matriz.
\]

\noindent\textbf{Comprobación de la respuesta} la aritmética del último paso reproduce la respuesta enunciada a la precisión con que está escrita. Es una reproducción aritmética; no valida el modelo físico ni sus hipótesis.

\noindent\textbf{Ejercicio 1 [PRB-0045]}

En Ada, el tipo \texttt{type Percent is range 0 .. 100;} comprende desde 0 hasta 100 ambos inclusive. ¿Cuántos valores distintos pertenecen al tipo (\texttt{Percent'Last - Percent'First + 1})?

\noindent\textbf{Respuesta} 101

\noindent\textbf{Ejercicio 2 [PRB-0046]}

En Fortran, una rebanada vectorial \texttt{u(2:n-1)} para \(n = 1002\) evalúa los puntos interiores de una malla 1D. ¿Cuántos puntos interiores actualiza la sentencia?

\noindent\textbf{Respuesta} 1000

\noindent\textbf{Ejercicio 3 [PRB-0047]}

Un programa Coarray Fortran se ejecuta en 16 imágenes (\texttt{num\_images() = 16}) y cada imagen aloja una malla local de \(50000\) flotantes \texttt{real64} (8 B). ¿Cuántos bytes suman las mallas de las 16 imágenes?

\noindent\textbf{Respuesta} 6400000

\noindent\textbf{Ejercicio 4 [PRB-0048]}

En un sistema de tiempo real Ada Ravenscar con periodo de marco mayor \(T = 100\text{ ms}\), tres tareas periódicas consumen \(15\text{ ms}\), \(22\text{ ms}\) y \(18\text{ ms}\) por cada \(100\text{ ms}\). ¿Cuál es el porcentaje entero de utilización total de CPU?

\noindent\textbf{Respuesta} 55

\subsection*{Cálculos reproducidos}

\(\left(8\cdot \left(\left(\left(4-1\right)\cdot 100\right)+\left(25-1\right)\right)\right) = 2592\)
\chapter{Auditoría de fuentes}

\noindent Estado de las fuentes: PENDING. 32 pendientes. Los conflictos almacenados quedan fuera (0). Las fuentes que la guía almacenada no cita quedan fuera. Recuento de no citadas: 0.

\noindent ¶1: dos testimonios.

\noindent ¶2: dos testimonios.

\noindent ¶3: dos testimonios.

\noindent ¶4: dos testimonios.

\noindent ¶5: dos testimonios.

\noindent ¶6: dos testimonios.

\noindent ¶7: dos testimonios.

\noindent ¶8: dos testimonios.

\noindent ¶9: dos testimonios.

\noindent ¶10: dos testimonios.

\noindent ¶11: dos testimonios.

\noindent ¶12: dos testimonios.

\noindent ¶13: dos testimonios.

\noindent ¶14: dos testimonios.

\noindent ¶15: dos testimonios.

\noindent ¶16: dos testimonios.

\noindent ¶17: dos testimonios.

\noindent ¶18: dos testimonios.

\noindent ¶19: dos testimonios.

\noindent ¶20: dos testimonios.

\noindent ¶21: dos testimonios.

\noindent ¶22: dos testimonios.

\noindent ¶23: dos testimonios.

\noindent ¶24: dos testimonios.

\noindent ¶25: dos testimonios.

\noindent ¶26: dos testimonios.

\noindent ¶27: dos testimonios.

\noindent ¶28: dos testimonios.

\noindent ¶29: dos testimonios.

\noindent ¶30: dos testimonios.

\noindent ¶31: dos testimonios.

\noindent ¶32: dos testimonios.

\noindent ¶33: dos testimonios.

\noindent ¶34: dos testimonios.

\noindent ¶35: dos testimonios.

\noindent ¶36: dos testimonios.

\noindent ¶37: dos testimonios.

\noindent ¶38: dos testimonios.

\noindent ¶39: dos testimonios.

\noindent ¶40: dos testimonios.

\noindent ¶41: dos testimonios.

\noindent ¶42: dos testimonios.

\noindent ¶43: dos testimonios.

\noindent ¶44: dos testimonios.

\noindent ¶45: dos testimonios.

\noindent ¶46: dos testimonios.

\noindent ¶47: dos testimonios.

\noindent ¶48: dos testimonios.

\medskip

\noindent Referencia Normativa src-iso-c17 (PENDING) --- Suplemento abierto: --- No es la bibliografía de la guía. --- \P1, \P2, \P3, \P4, \P5, \P6, \P7, \P8, \P9, \P10, \P11, \P12, \P13, \P14, \P15, \P16, \P17, \P18, \P31, \P32, \P33, \P34, \P35, \P36.

\noindent Referencia Normativa src-bryant-ohallaron (PENDING) --- Suplemento abierto: --- No es la bibliografía de la guía. --- \P1, \P2, \P3, \P4, \P5, \P6, \P7, \P8, \P9, \P10, \P11, \P12, \P19, \P20, \P21, \P22, \P23, \P24, \P31, \P32, \P33, \P34, \P35, \P36, \P37, \P38, \P39, \P40, \P41, \P42.

\noindent Referencia Normativa src-posix-2017 (PENDING) --- Suplemento abierto: --- No es la bibliografía de la guía. --- \P13, \P14, \P15, \P16, \P17, \P18.

\noindent Referencia Normativa src-iso-cpp20 (PENDING) --- Suplemento abierto: --- No es la bibliografía de la guía. --- \P19, \P20, \P21, \P22, \P23, \P24, \P25, \P26, \P27, \P28, \P29, \P30.

\noindent Referencia Normativa src-cormen-clrs (PENDING) --- Suplemento abierto: --- No es la bibliografía de la guía. --- \P25, \P26, \P27, \P28, \P29, \P30.

\noindent Referencia Normativa src-patterson-hennessy (PENDING) --- Suplemento abierto: --- No es la bibliografía de la guía. --- \P37, \P38, \P39, \P40, \P41, \P42.

\noindent IEEE Standard for Floating-Point Arithmetic (IEEE 754-2019) / Wikipedia Summary (PENDING) --- Suplemento abierto: --- No es la bibliografía de la guía. --- \P43, \P44, \P45, \P46, \P47, \P48.

\noindent Referencia Normativa src-knuth-taocp (PENDING) --- Suplemento abierto: --- No es la bibliografía de la guía. --- \P43, \P44, \P45, \P46, \P47, \P48.

\noindent REV-0001: cap01-c-fundamentos-modelo-compilacion-tipos-principiante scientific

\noindent REV-0002: cap02-c-punteros-memoria-dinamica-estructuras-aliasing scientific

\noindent REV-0003: cap03-c-experto-atomicos-posix-sanitizers-ecosistema-proyectos scientific

\noindent REV-0004: cap04-cpp-fundamentos-raii-move-semantics-smart-pointers scientific

\noindent REV-0005: cap05-cpp-templates-concepts-ranges-concurrencia-proyectos scientific

\noindent REV-0006: cap06-monografia-zig-asignadores-comptime-interoperabilidad-c scientific

\noindent REV-0007: cap07-monografia-ensamblador-arm64-x86-64-abi-simd scientific

\noindent REV-0008: cap08-monografias-ada-spark-tiempo-real-y-fortran-cientifico scientific
\chapter{Plan de estudio}

\section*{Monografía Autónoma de C (ISO/IEC 9899:2018 / C23)}

\subsection*{\checkmark Monografía de C (I \(-\) Bloques A, B y C): Historia, modelo de traducción ISO C17/C23, toolchain y nivel principiante}
\begin{itemize}
\item Comprender la evolución histórica de C desde UNIX en Bell Labs (Dennis Ritchie, 1972) hasta los estándares ANSI C89, C99, C11, C17 y C23.
\item Dominar las 8 fases de traducción del estándar ISO C, el preprocesador, la compilación separada por unidades de traducción y el enlazado ELF/Mach-O.
\item Utilizar con precisión los tipos enteros de anchura exacta (<stdint.h>), promociones enteras, operadores bit a bit y estructuras de control.
\end{itemize}
\noindent 2 conceptos planificados.

\subsection*{\checkmark Monografía de C (II \(-\) Bloques D y E): Punteros, aritmética de direcciones, memoria dinámica, structs y strict aliasing}
\begin{itemize}
\item Dominar la semántica de punteros, el decaimiento de arrays a punteros, punteros a funciones y la aritmética escalada de direcciones.
\item Gestionar el ciclo de vida en el montículo con malloc, calloc, realloc y free, y construir asignadores lineales de región (Arena Allocators).
\item Aplicar la regla de alias estricto (strict aliasing), el calificador restrict y la lectura binaria segura mediante memcpy.
\end{itemize}
\noindent 2 conceptos planificados.

\subsection*{\checkmark Monografía de C (III \(-\) Bloques F, G y H): Concurrencia C11/POSIX, Sanitizers, Ecosistema y Proyectos}
\begin{itemize}
\item Programar concurrencia segura en C mediante <stdatomic.h>, pthreads, descriptores de fichero POSIX y mmap.
\item Integrar AddressSanitizer (ASan), UndefinedBehaviorSanitizer (UBSan), ThreadSanitizer (TSan) y analizadores estáticos en CMake/Meson.
\item Auditar y extender el proyecto completo `projects/c-arena-allocator` junto con los 4 proyectos integradores en C.
\end{itemize}
\noindent 2 conceptos planificados.

\section*{Monografía Autónoma de C++ (ISO/IEC 14882:2020 / C++23)}

\subsection*{\checkmark Monografía de C++ (I \(-\) Bloques A\(-\)D): De C++98 a C++20/23, RAII, Semántica de Movimiento y Punteros Inteligentes}
\begin{itemize}
\item Comprender la filosofía de C++ moderno (Bjarne Stroustrup): tipos definidos por el usuario con coste cero y determinismo RAII.
\item Aplicar la Regla de Cero, la Regla de Cinco, categorías de valores (lvalue, prvalue, xvalue) y semántica de movimiento (std::move).
\item Gestionar recursos dinámicos sin fugas mediante std::unique\_ptr, std::shared\_ptr, std::weak\_ptr, std::span y std::string\_view.
\end{itemize}
\noindent 2 conceptos planificados.

\subsection*{\checkmark Monografía de C++ (II \(-\) Bloques E\(-\)H): Plantillas, Concepts C++20, std::ranges, Concurrencia, Corrutinas y Proyectos}
\begin{itemize}
\item Escribir código genérico restringido semánticamente mediante C++20 Concepts (`requires`), plantillas variádicas y `if constexpr`.
\item Componer tuberías perezosas de vistas con `std::ranges` y `std::views` evitando copias intermedias.
\item Dominar la concurrencia en C++20 (`std::jthread`, `std::stop\_token`, `std::atomic`, `std::latch`), el ecosistema CMake/Conan/vcpkg y el proyecto `projects/cpp-task-scheduler`.
\end{itemize}
\noindent 2 conceptos planificados.

\section*{Monografías Especializadas de Sistemas: Zig, Ensamblador, Ada y Fortran}

\subsection*{\checkmark Monografía de Zig (Bloques A\(-\)H): Asignadores explícitos, metaprogramación comptime, uniones de error y toolchain C}
\begin{itemize}
\item Comprender los principios de diseño de Zig: ausencia de flujo de control oculto, ausencia de asignaciones implícitas de montículo y ausencia de preprocesador.
\item Programar con uniones de error (`!T`), limpieza diferida determinista (`defer` y `errdefer`), rebanadas con longitud y asignadores explícitos (`std.mem.Allocator`).
\item Utilizar `comptime` para tipos genéricos y evaluación en tiempo de compilación, e integrar cabeceras C con `@cImport`.
\end{itemize}
\noindent 2 conceptos planificados.

\subsection*{\checkmark Monografía de Ensamblador ARM64 (AArch64) y x86-64 (Bloques A\(-\)H): Registros, ABI, Pila y SIMD}
\begin{itemize}
\item Distinguir la arquitectura RISC de carga/almacenamiento AArch64 (ARM64) frente a la arquitectura CISC de longitud variable x86-64.
\item Dominar las convenciones de llamada ABI System V AMD64 (`rdi, rsi, rdx, rcx, r8, r9` y retorno en `rax`) y AAPCS64 (`x0\(-\)x7` y retorno en `x0`).
\item Analizar el prólogo y epílogo de marcos de pila, el alineamiento obligatorio de 16 bytes del puntero de pila (`sp`/`rsp`) e instrucciones vectoriales SIMD.
\end{itemize}
\noindent 2 conceptos planificados.

\subsection*{\checkmark Monografías de Ada/SPARK (Sistemas Críticos) y Fortran Moderno (Cómputo Científico HPC) (Bloques A\(-\)H)}
\begin{itemize}
\item Diseñar tipos escalares con rangos restringidos, contratos ejecutables (`Pre`, `Post`, `Type\_Invariant`) y tareas críticas bajo el perfil Ravenscar en Ada 2022 y SPARK.
\item Comprender cómo el demostrador formal GNATprove en SPARK demuestra estáticamente la ausencia de errores en tiempo de ejecución (AoRTE).
\item Programar cálculo numérico de alto rendimiento en Fortran 2018/2023 mediante operaciones matriciales intrínsecas, `intent(in/out)`, módulos y paralelismo Coarray.
\end{itemize}
\noindent 2 conceptos planificados.
\backmatter
\begin{thebibliography}{99}
\bibitem{src1} Alan M. Turing \(-\) On Computable Numbers, with an Application to the Entscheidungsproblem (1936). \url{https://www.cs.virginia.edu/~robins/Turing_Paper_1936.pdf}
\bibitem{src2} John von Neumann \(-\) First Draft of a Report on the EDVAC (1945). \url{https://web.mit.edu/STS.035/www/PDFs/edvac.pdf}
\bibitem{src3} IEEE Standard for Floating-Point Arithmetic (IEEE 754-2019) / Wikipedia Summary. \url{https://en.wikipedia.org/wiki/IEEE_754}
\bibitem{src4} RFC 3629 \(-\) UTF-8, a transformation format of ISO 10646 (IETF Standard). \url{https://datatracker.ietf.org/doc/html/rfc3629}
\bibitem{src5} ACM / IEEE-CS / AAAI Computer Science Curricula 2023 (CS2023). \url{https://csed.acm.org/}
\bibitem{src6} MIT OpenCourseWare 6.006 \(-\) Introduction to Algorithms. \url{https://ocw.mit.edu/courses/6-006-introduction-to-algorithms-spring-2020/}
\bibitem{src7} The Rust Programming Language (Klabnik \& Nichols \(-\) Official Rust Documentation). \url{https://doc.rust-lang.org/book/}
\bibitem{src8} The Rust Reference \(-\) Memory Model, Type System, Lifetimes, and Unsafe Rust. \url{https://doc.rust-lang.org/reference/}
\bibitem{src9} The Rustonomicon \(-\) The Dark Arts of Unsafe Rust. \url{https://doc.rust-lang.org/nomicon/}
\bibitem{src10} C Programming Language Reference (ISO/IEC 9899 C17/C23 \(-\) cppreference). \url{https://en.cppreference.com/w/c}
\bibitem{src11} C++ Reference (ISO/IEC 14882 C++20/C++23 \(-\) cppreference). \url{https://en.cppreference.com/w/cpp}
\bibitem{src12} The Python Language Reference \& Standard Library Documentation (Python 3.12\(-\)3.14). \url{https://docs.python.org/3/reference/index.html}
\bibitem{src13} ECMA-262 \(-\) ECMAScript 2024 Language Specification. \url{https://tc39.es/ecma262/}
\bibitem{src14} The TypeScript Handbook \(-\) Structural Typing, Narrowing, and Type Manipulation. \url{https://www.typescriptlang.org/docs/handbook/intro.html}
\bibitem{src15} Microsoft .NET \& C\# Language Specification (ECMA-334 / .NET 8 \& 9). \url{https://learn.microsoft.com/en-us/dotnet/csharp/}
\bibitem{src16} The Java Language Specification \& JVM Specification (Java SE 21 LTS). \url{https://docs.oracle.com/javase/specs/jls/se21/html/index.html}
\bibitem{src17} The Go Programming Language Specification \& Memory Model. \url{https://go.dev/ref/spec}
\bibitem{src18} Kotlin Language Documentation \& Coroutines Guide (Kotlin 2.0). \url{https://kotlinlang.org/docs/home.html}
\bibitem{src19} The Swift Programming Language (Swift 6.0 Official Book). \url{https://docs.swift.org/swift-book/documentation/the-swift-programming-language/}
\bibitem{src20} E. F. Codd \(-\) A Relational Model of Data for Large Shared Data Banks (1970) \& PostgreSQL Docs. \url{https://www.postgresql.org/docs/current/}
\bibitem{src21} SQLite 3.54 Transactional SQL Database Engine Documentation. \url{https://www.sqlite.org/docs.html}
\bibitem{src22} GNU Bash Reference Manual \& IEEE Std 1003.1 POSIX Shell Specification. \url{https://www.gnu.org/software/bash/manual/bash.html}
\bibitem{src23} WHATWG HTML Living Standard \& W3C CSS Specifications. \url{https://html.spec.whatwg.org/multipage/}
\bibitem{src24} RFC 8259 \(-\) The JavaScript Object Notation (JSON) Data Interchange Format \& YAML 1.2.2. \url{https://datatracker.ietf.org/doc/html/rfc8259}
\bibitem{src25} IEEE Computer Society \(-\) Guide to the Software Engineering Body of Knowledge (SWEBOK v4.0). \url{https://www.computer.org/education/bodies-of-knowledge/software-engineering}
\bibitem{src26} Referencia Normativa src-iso-c17. \url{https://doc.rust-lang.org/reference/src-iso-c17.html}
\bibitem{src27} Referencia Normativa src-bryant-ohallaron. \url{https://doc.rust-lang.org/reference/src-bryant-ohallaron.html}
\bibitem{src28} Referencia Normativa src-posix-2017. \url{https://doc.rust-lang.org/reference/src-posix-2017.html}
\bibitem{src29} Referencia Normativa src-iso-cpp20. \url{https://doc.rust-lang.org/reference/src-iso-cpp20.html}
\bibitem{src30} Referencia Normativa src-cormen-clrs. \url{https://doc.rust-lang.org/reference/src-cormen-clrs.html}
\bibitem{src31} Referencia Normativa src-patterson-hennessy. \url{https://doc.rust-lang.org/reference/src-patterson-hennessy.html}
\bibitem{src32} Referencia Normativa src-knuth-taocp. \url{https://doc.rust-lang.org/reference/src-knuth-taocp.html}
\end{thebibliography}
\end{document}
